% Options for packages loaded elsewhere
\PassOptionsToPackage{unicode}{hyperref}
\PassOptionsToPackage{hyphens}{url}
%
\documentclass[
  9pt,
]{article}
\usepackage{amsmath,amssymb}
\usepackage{setspace}
\usepackage{iftex}
\ifPDFTeX
  \usepackage[T1]{fontenc}
  \usepackage[utf8]{inputenc}
  \DeclareUnicodeCharacter{2192}{\ensuremath{\rightarrow}}
  \DeclareUnicodeCharacter{2194}{\ensuremath{\leftrightarrow}}
  \DeclareUnicodeCharacter{21C4}{\ensuremath{\rightleftarrows}}
  \DeclareUnicodeCharacter{2198}{\ensuremath{\searrow}}
  \DeclareUnicodeCharacter{2265}{\ensuremath{\geq}}
  \DeclareUnicodeCharacter{2264}{\ensuremath{\leq}}
  \DeclareUnicodeCharacter{00D7}{\ensuremath{\times}}
  \usepackage{textcomp} % provide euro and other symbols
\else % if luatex or xetex
  \usepackage{unicode-math} % this also loads fontspec
  \defaultfontfeatures{Scale=MatchLowercase}
  \defaultfontfeatures[\rmfamily]{Ligatures=TeX,Scale=1}
\fi
\usepackage{lmodern}
\ifPDFTeX\else
  % xetex/luatex font selection
  \setmainfont[]{DejaVu Sans}
  \setmonofont[]{DejaVu Sans Mono}
\fi
% Use upquote if available, for straight quotes in verbatim environments
\IfFileExists{upquote.sty}{\usepackage{upquote}}{}
\IfFileExists{microtype.sty}{% use microtype if available
  \usepackage[]{microtype}
  \UseMicrotypeSet[protrusion]{basicmath} % disable protrusion for tt fonts
}{}
\makeatletter
\@ifundefined{KOMAClassName}{% if non-KOMA class
  \IfFileExists{parskip.sty}{%
    \usepackage{parskip}
  }{% else
    \setlength{\parindent}{0pt}
    \setlength{\parskip}{6pt plus 2pt minus 1pt}}
}{% if KOMA class
  \KOMAoptions{parskip=half}}
\makeatother
\usepackage{xcolor}
\usepackage[margin=0.74in]{geometry}
\usepackage{color}
\usepackage{fancyvrb}
\newcommand{\VerbBar}{|}
\newcommand{\VERB}{\Verb[commandchars=\\\{\}]}
\DefineVerbatimEnvironment{Highlighting}{Verbatim}{commandchars=\\\{\}}
% Add ',fontsize=\small' for more characters per line
\usepackage{framed}
\definecolor{shadecolor}{RGB}{248,248,248}
\newenvironment{Shaded}{\begin{snugshade}}{\end{snugshade}}
\newcommand{\AlertTok}[1]{\textcolor[rgb]{0.94,0.16,0.16}{#1}}
\newcommand{\AnnotationTok}[1]{\textcolor[rgb]{0.56,0.35,0.01}{\textbf{\textit{#1}}}}
\newcommand{\AttributeTok}[1]{\textcolor[rgb]{0.13,0.29,0.53}{#1}}
\newcommand{\BaseNTok}[1]{\textcolor[rgb]{0.00,0.00,0.81}{#1}}
\newcommand{\BuiltInTok}[1]{#1}
\newcommand{\CharTok}[1]{\textcolor[rgb]{0.31,0.60,0.02}{#1}}
\newcommand{\CommentTok}[1]{\textcolor[rgb]{0.56,0.35,0.01}{\textit{#1}}}
\newcommand{\CommentVarTok}[1]{\textcolor[rgb]{0.56,0.35,0.01}{\textbf{\textit{#1}}}}
\newcommand{\ConstantTok}[1]{\textcolor[rgb]{0.56,0.35,0.01}{#1}}
\newcommand{\ControlFlowTok}[1]{\textcolor[rgb]{0.13,0.29,0.53}{\textbf{#1}}}
\newcommand{\DataTypeTok}[1]{\textcolor[rgb]{0.13,0.29,0.53}{#1}}
\newcommand{\DecValTok}[1]{\textcolor[rgb]{0.00,0.00,0.81}{#1}}
\newcommand{\DocumentationTok}[1]{\textcolor[rgb]{0.56,0.35,0.01}{\textbf{\textit{#1}}}}
\newcommand{\ErrorTok}[1]{\textcolor[rgb]{0.64,0.00,0.00}{\textbf{#1}}}
\newcommand{\ExtensionTok}[1]{#1}
\newcommand{\FloatTok}[1]{\textcolor[rgb]{0.00,0.00,0.81}{#1}}
\newcommand{\FunctionTok}[1]{\textcolor[rgb]{0.13,0.29,0.53}{\textbf{#1}}}
\newcommand{\ImportTok}[1]{#1}
\newcommand{\InformationTok}[1]{\textcolor[rgb]{0.56,0.35,0.01}{\textbf{\textit{#1}}}}
\newcommand{\KeywordTok}[1]{\textcolor[rgb]{0.13,0.29,0.53}{\textbf{#1}}}
\newcommand{\NormalTok}[1]{#1}
\newcommand{\OperatorTok}[1]{\textcolor[rgb]{0.81,0.36,0.00}{\textbf{#1}}}
\newcommand{\OtherTok}[1]{\textcolor[rgb]{0.56,0.35,0.01}{#1}}
\newcommand{\PreprocessorTok}[1]{\textcolor[rgb]{0.56,0.35,0.01}{\textit{#1}}}
\newcommand{\RegionMarkerTok}[1]{#1}
\newcommand{\SpecialCharTok}[1]{\textcolor[rgb]{0.81,0.36,0.00}{\textbf{#1}}}
\newcommand{\SpecialStringTok}[1]{\textcolor[rgb]{0.31,0.60,0.02}{#1}}
\newcommand{\StringTok}[1]{\textcolor[rgb]{0.31,0.60,0.02}{#1}}
\newcommand{\VariableTok}[1]{\textcolor[rgb]{0.00,0.00,0.00}{#1}}
\newcommand{\VerbatimStringTok}[1]{\textcolor[rgb]{0.31,0.60,0.02}{#1}}
\newcommand{\WarningTok}[1]{\textcolor[rgb]{0.56,0.35,0.01}{\textbf{\textit{#1}}}}
\usepackage{longtable,booktabs,array}
\usepackage{calc} % for calculating minipage widths
% Correct order of tables after \paragraph or \subparagraph
\usepackage{etoolbox}
\makeatletter
\patchcmd\longtable{\par}{\if@noskipsec\mbox{}\fi\par}{}{}
\makeatother
% Allow footnotes in longtable head/foot
\IfFileExists{footnotehyper.sty}{\usepackage{footnotehyper}}{\usepackage{footnote}}
\makesavenoteenv{longtable}
\usepackage{graphicx}
\makeatletter
\def\maxwidth{\ifdim\Gin@nat@width>\linewidth\linewidth\else\Gin@nat@width\fi}
\def\maxheight{\ifdim\Gin@nat@height>\textheight\textheight\else\Gin@nat@height\fi}
\makeatother
% Scale images if necessary, so that they will not overflow the page
% margins by default, and it is still possible to overwrite the defaults
% using explicit options in \includegraphics[width, height, ...]{}
\setkeys{Gin}{width=\maxwidth,height=\maxheight,keepaspectratio}
% Set default figure placement to htbp
\makeatletter
\def\fps@figure{htbp}
\makeatother
\setlength{\emergencystretch}{3em} % prevent overfull lines
\providecommand{\tightlist}{%
  \setlength{\itemsep}{0pt}\setlength{\parskip}{0pt}}
\setcounter{secnumdepth}{-\maxdimen} % remove section numbering
\usepackage{microtype}
\usepackage{xcolor}
\definecolor{navyblue}{HTML}{18334A}
\definecolor{steelblue}{HTML}{275874}
\definecolor{mutedgray}{HTML}{62707B}
\usepackage{fancyhdr}
\pagestyle{fancy}
\fancyhf{}
\fancyhead[L]{\small\textcolor{mutedgray}{AI Agent Engineering | Curriculum v1.0}}
\fancyhead[R]{\small\textcolor{mutedgray}{November 2026}}
\fancyfoot[C]{\small\textcolor{mutedgray}{\thepage}}
\renewcommand{\headrulewidth}{0.25pt}
\setlength{\headheight}{14pt}
\usepackage{titlesec}
\titleformat{\section}{\Large\bfseries\color{navyblue}}{\thesection}{0.65em}{}
\titleformat{\subsection}{\large\bfseries\color{steelblue}}{\thesubsection}{0.65em}{}
\titleformat{\subsubsection}{\normalsize\bfseries\color{steelblue}}{\thesubsubsection}{0.65em}{}
\usepackage{enumitem}
\setlist[itemize]{topsep=2pt,itemsep=1pt,parsep=1pt}
\setlist[enumerate]{topsep=2pt,itemsep=1pt,parsep=1pt}
\AtBeginDocument{\hypersetup{colorlinks=true,linkcolor=steelblue,urlcolor=steelblue,citecolor=steelblue}}
\usepackage{float}
\usepackage{needspace}
\clubpenalty=10000
\widowpenalty=10000
\emergencystretch=3em
\usepackage{pdflscape}
\ifLuaTeX
  \usepackage{selnolig}  % disable illegal ligatures
\fi
\usepackage{bookmark}
\IfFileExists{xurl.sty}{\usepackage{xurl}}{} % add URL line breaks if available
\urlstyle{same}
\hypersetup{
  pdftitle={AI Agent Engineering},
  hidelinks,
  pdfcreator={LaTeX via pandoc}}

\title{AI Agent Engineering}
\usepackage{etoolbox}
\makeatletter
\providecommand{\subtitle}[1]{% add subtitle to \maketitle
  \apptocmd{\@title}{\par {\large #1 \par}}{}{}
}
\makeatother
\subtitle{Comprehensive Developer Curriculum: Frameworks, Harnesses, MCP
and Enterprise Security}
\author{}
\date{November 12, 2026}

\begin{document}
\maketitle

{
\setcounter{tocdepth}{2}
\tableofcontents
}
\setstretch{1.08}
\section{AI Agent Engineering: Agent Frameworks, Harness Architecture,
Model Context Protocol, and Enterprise-Grade Autonomous
Systems}\label{ai-agent-engineering-agent-frameworks-harness-architecture-model-context-protocol-and-enterprise-grade-autonomous-systems}

\textbf{Comprehensive Developer Curriculum and Engineering Syllabus}\\
\textbf{Edition:} 1.0 · \textbf{Source review:} November 12, 2026\\
\textbf{Duration:} 16 weeks · \textbf{Workload:} 10 hours/week (160
total) · \textbf{Level:} Intermediate to advanced\\
\textbf{Primary language:} Python · \textbf{Optional tracks:} TypeScript
and Java · \textbf{Capstone:} Governed GitHub CodeQL investigation agent

\begin{quote}
\textbf{Provenance.} This is a new instructional design, \textbf{not} a
transcription of the supplied architecture reference. The reference
document establishes the responsibilities of the model, framework,
harness, MCP, and application governance; the learning sequence,
durations, acceptance criteria, test thresholds, and capstone
assignments below are proposed educational decisions. Course URLs and
externally sourced facts are recorded in sections 9--11 and 23.
\end{quote}

\newpage
\subsection{1. Curriculum Executive
Summary}\label{curriculum-executive-summary}

\subsubsection{Mission and scope}\label{mission-and-scope}

The program prepares experienced developers to build AI agent systems
with independently testable runtime, integration, security, and
business-control boundaries. Learners begin with a small framework-free
runner, compare it to a production framework, construct a governed
harness, add MCP tools, implement state and retrieval, and finish with
an AppSec investigation agent. All remote interactions are optional
during the foundational phase; mock providers make the core exercises
runnable without a paid model API or GitHub organization access.

\textbf{Nonnegotiable distinction:} the model proposes a tool call; the
harness validates and dispatches it; the integration layer carries the
request; the downstream resource enforces its own permissions;
application policy or an authorized human decides whether a sensitive
action is permitted. A framework can supply components of the harness
but does not automatically solve authorization or auditability.

\subsubsection{Instructional outcomes and
portfolio}\label{instructional-outcomes-and-portfolio}

\begin{itemize}
\tightlist
\item
  \textbf{16} structured weekly units and \textbf{16} engineering labs,
  one per week.
\item
  \textbf{12} technical modules spanning foundations through deployment
  and governance.
\item
  \textbf{Two stage gates:} W8 midterm and W16 capstone
  production-readiness review.
\item
  \textbf{Three parallel evidence streams:} runnable code and tests;
  architecture/security documentation; measurable quality/operational
  evidence.
\item
  \textbf{One representative final system:} GitHub CodeQL alert
  investigation and recommendations, \textbf{read-only by default}.
  Dismissal actions are demonstrated only against mocks unless
  separately authorized outside the course.
\end{itemize}

\begin{longtable}[]{@{}
  >{\raggedright\arraybackslash}p{(\columnwidth - 6\tabcolsep) * \real{0.2143}}
  >{\raggedleft\arraybackslash}p{(\columnwidth - 6\tabcolsep) * \real{0.2857}}
  >{\raggedleft\arraybackslash}p{(\columnwidth - 6\tabcolsep) * \real{0.2857}}
  >{\raggedright\arraybackslash}p{(\columnwidth - 6\tabcolsep) * \real{0.2143}}@{}}
\toprule\noalign{}
\begin{minipage}[b]{\linewidth}\raggedright
Work category
\end{minipage} & \begin{minipage}[b]{\linewidth}\raggedleft
Hours
\end{minipage} & \begin{minipage}[b]{\linewidth}\raggedleft
Share
\end{minipage} & \begin{minipage}[b]{\linewidth}\raggedright
Evidence
\end{minipage} \\
\midrule\noalign{}
\endhead
\bottomrule\noalign{}
\endlastfoot
Lectures, demonstrations, official tutorials & 36 & 22.5\% & Notes and
short quizzes \\
Technical readings and source comparison & 24 & 15.0\% & Architecture
annotations \\
Programming and laboratory implementation & 64 & 40.0\% & Reviewed
commits and test runs \\
Testing, security exercises and evaluations & 20 & 12.5\% & CI reports,
test and policy evidence \\
Reviews, reports, midterm and capstone handoff & 16 & 10.0\% & ADRs,
final readiness packet \\
\textbf{Total} & \textbf{160} & \textbf{100\%} & \\
\end{longtable}

\textbf{Pacing rule:} each week budgets 2 h lecture/tutorial, 1.5 h
reading, 4 h coding, 1.5 h testing, and 1 h evidence/review; weeks 8, 15
and 16 reallocate time toward demonstrations and capstone integration
while remaining within 10 h.

\subsection{2. Course Description and Learning
Philosophy}\label{course-description-and-learning-philosophy}

This is a systems engineering course, not a chatbot prompt-writing
class. Instruction uses four repeated activities: \textbf{explain a
boundary → implement the smallest enforceable mechanism → attack and
test the mechanism → preserve reproducible evidence}. Code runs behind
explicit typed interfaces and fake model/GitHub adapters before any live
deployment. Models may be nondeterministic; authorization, input
validation, resource limits, evidence contracts and deployment decisions
must not be.

The supplied architecture reference describes a bounded lifecycle
(intake, initialization, proposal, validation, dispatch, observation,
approval, verification and closure). Each laboratory adds one capability
to that lifecycle. Its stable exit states are \texttt{completed},
\texttt{needs\_approval}, \texttt{blocked}, \texttt{failed},
\texttt{cancelled}, and \texttt{budget\_exhausted}. A \texttt{completed}
result requires independently defined success predicates, not merely a
fluent model answer.

\textbf{Design constraints:} local-first; minimal secrets exposure;
testable ports and adapters; no production dismissals; deterministic
policy gates; full source revision attribution; reproducible environment
manifests; explicit supported MCP version per tool integration. Never
imply that the presence of an MCP tool authorizes its use.

\subsection{3. Target Audience and Prerequisite
Analysis}\label{target-audience-and-prerequisite-analysis}

\textbf{Roles:} backend/platform software engineer; AppSec or DevSecOps
engineer; reliability engineer; technical architect. The baseline
assumes Git, Linux, REST/JSON, data structures, modular software design,
code review, test assertions, and introductory CI/CD. No formal AI/ML
mathematics course is needed to begin.

\begin{longtable}[]{@{}
  >{\raggedright\arraybackslash}p{(\columnwidth - 4\tabcolsep) * \real{0.3333}}
  >{\raggedright\arraybackslash}p{(\columnwidth - 4\tabcolsep) * \real{0.3333}}
  >{\raggedright\arraybackslash}p{(\columnwidth - 4\tabcolsep) * \real{0.3333}}@{}}
\toprule\noalign{}
\begin{minipage}[b]{\linewidth}\raggedright
Area
\end{minipage} & \begin{minipage}[b]{\linewidth}\raggedright
Prerequisite status
\end{minipage} & \begin{minipage}[b]{\linewidth}\raggedright
Working proficiency target and bridge
\end{minipage} \\
\midrule\noalign{}
\endhead
\bottomrule\noalign{}
\endlastfoot
Python 3 typing, \texttt{dataclass}, \texttt{Protocol}, exceptions,
pytest & \textbf{Essential by W1} & Implement typed DTOs, test mocks,
and exceptions; 4-hour remedial bridge \\
\texttt{async}/\texttt{await}, timeouts, cancellation &
\textbf{Essential by W2} & Compose coroutines without leaking tasks;
3-hour bridge \\
HTTP status codes, JSON Schema, auth headers & \textbf{Essential by W1}
& Understand schemas, pagination, 401/403/429; assumed baseline \\
SQLite and transactions & \textbf{By W4} & Store checkpoints and reject
stale updates; concurrent bridge \\
OAuth 2.0, GitHub Apps and token scopes & \textbf{By W7} & Separate
authentication from authorization; concurrent 3-hour bridge \\
Docker, container limits and Linux permissions & \textbf{By W5/W14} &
Run as non-root with restricted network/filesystem; concurrent bridge \\
Retrieval, information retrieval measures & \textbf{By W8} & Interpret
BM25, embeddings, precision@k; taught in course \\
Threat modeling and identity controls & \textbf{By W9} & STRIDE-like
abuse cases and least privilege; concurrent bridge \\
TypeScript/Node or Java/Spring Boot & \textbf{Optional advanced track} &
Reimplement typed MCP client/server interfaces after W7 \\
AWS, IaC, messaging/event buses & \textbf{Optional advanced deployment}
& Compare Lambda/Fargate/containers without cloud spend \\
\end{longtable}

\textbf{Entry diagnostic (30 points; 45 minutes):} 5 questions Python
types/async; 5 API/JSON/schema; 5 software tests; 5 Git/CI; 5
trust/auth; 5 LLM/tool concepts. One point per correct answer.
\textbf{Ready:} ≥24 and no essential domain below 3/5.
\textbf{Bridge-first:} 17--23 or any essential domain below 3/5.
\textbf{Foundation preparation:} ≤16. A sample practical diagnostic asks
candidates to validate an untrusted tool JSON object, explain 401 versus
403, design a timeout test, identify a non-idempotent retry, and state
why a model's instruction to dismiss an alert is not authorization.

\textbf{Bridge assignment:} submit \texttt{prereq\_check.py} with an
asynchronous \texttt{fetch\_fixture()} wrapped in a timeout, a typed
schema validator and a five-case pytest suite; no API key required.

\subsection{4. Curriculum Goals and Measurable Learning
Outcomes}\label{curriculum-goals-and-measurable-learning-outcomes}

\begin{longtable}[]{@{}
  >{\raggedright\arraybackslash}p{(\columnwidth - 4\tabcolsep) * \real{0.3333}}
  >{\raggedright\arraybackslash}p{(\columnwidth - 4\tabcolsep) * \real{0.3333}}
  >{\raggedright\arraybackslash}p{(\columnwidth - 4\tabcolsep) * \real{0.3333}}@{}}
\toprule\noalign{}
\begin{minipage}[b]{\linewidth}\raggedright
ID
\end{minipage} & \begin{minipage}[b]{\linewidth}\raggedright
By graduation, the learner can\ldots{}
\end{minipage} & \begin{minipage}[b]{\linewidth}\raggedright
Primary measurement
\end{minipage} \\
\midrule\noalign{}
\endhead
\bottomrule\noalign{}
\endlastfoot
LO01 & \textbf{Classify} the five architecture layers and justify their
decision rights & W1 boundary diagram and oral defense \\
LO02 & \textbf{Implement} a bounded framework-free model/tool loop & W2
finite-state tests and run trace \\
LO03 & \textbf{Compare and apply} framework orchestration and durable
checkpoints & W3--4 parity tests, ADR \\
LO04 & \textbf{Implement} harness policy checks, budgets, compaction and
recovery & W5 sandbox tests; W4 checkpoint replay \\
LO05 & \textbf{Build and verify} pinned MCP server/client
interoperability & W6--7 protocol tests \\
LO06 & \textbf{Evaluate} retrieval and provenance for repository
evidence & W8 benchmark and evidence ledger \\
LO07 & \textbf{Threat-model and mitigate} indirect prompt injection and
excessive agency & W9 abuse-case suite \\
LO08 & \textbf{Enforce} independent policy decisions and approvals & W10
negative authorization tests \\
LO09 & \textbf{Orchestrate} scoped multi-agent delegation and failure
handling & W11 worker isolation/recovery suite \\
LO10 & \textbf{Measure} agent quality with regression datasets and
metrics & W12 reproducible evaluation \\
LO11 & \textbf{Instrument and operate} agent services using traces and
SLOs & W13 observability and fault injection \\
LO12 & \textbf{Package and release} secure agent workloads through CI/CD
& W14 verified build and release checklist \\
LO13 & \textbf{Integrate} the five layers in read-only AppSec operations
& W15 CodeQL evidence-backed recommendations \\
LO14 & \textbf{Defend} a release-readiness decision using independent
evidence & W16 architecture and security panel \\
\end{longtable}

\subsection{5. Prerequisite Dependency
Graph}\label{prerequisite-dependency-graph}

\begin{figure}
\centering
\fbox{\begin{minipage}{0.91\linewidth}
\centering\small
\textbf{Python typing + asynchronous programming}\quad
\textbf{REST + JSON Schema}\\[4pt]
$\Downarrow$\\[3pt]
\textbf{Weeks 1--2: Model boundaries and framework-free runner}\\[3pt]
$\Downarrow$\\[3pt]
\textbf{Weeks 3--5: Agent framework and harness runtime}\\[3pt]
$\Downarrow$\\[3pt]
\textbf{Weeks 6--7: MCP hosts, clients, and servers}\\[3pt]
$\Downarrow$\\[3pt]
\textbf{Weeks 8--14: Context, security, evaluation, and deployment}\\[3pt]
$\Downarrow$\\[3pt]
\textbf{Weeks 15--16: Enterprise GitHub AppSec capstone}\\[5pt]
\textit{Optional parallel specializations: TypeScript, Java, AWS}
\end{minipage}}
\caption{Prerequisite and progression graph}
\end{figure}

The diagram distinguishes required conceptual dependencies (solid
arrows) from optional implementation specialization (dashed arrows).
\texttt{Python\ typing\ +\ async} and \texttt{REST\ +\ JSON\ Schema}
precede the custom harness;
\texttt{boundaries\ →\ graph\ orchestration\ →\ durable\ state} precedes
MCP/GitHub; evidence ingestion precedes threat-modeling, evaluation and
the final capstone. TypeScript and Java are optional integrations, not
prerequisites to finish the Python core.

\subsection{6. Complete 16-Week Instructional
Roadmap}\label{complete-16-week-instructional-roadmap}

\begin{landscape}
\begin{longtable}[]{@{}
  >{\raggedright\arraybackslash}p{(\columnwidth - 10\tabcolsep) * \real{0.0700}}
  >{\raggedright\arraybackslash}p{(\columnwidth - 10\tabcolsep) * \real{0.0800}}
  >{\raggedright\arraybackslash}p{(\columnwidth - 10\tabcolsep) * \real{0.4900}}
  >{\raggedright\arraybackslash}p{(\columnwidth - 10\tabcolsep) * \real{0.0800}}
  >{\raggedright\arraybackslash}p{(\columnwidth - 10\tabcolsep) * \real{0.2000}}
  >{\raggedleft\arraybackslash}p{(\columnwidth - 10\tabcolsep) * \real{0.0800}}@{}}
\toprule\noalign{}
\begin{minipage}[b]{\linewidth}\raggedright
Week
\end{minipage} & \begin{minipage}[b]{\linewidth}\raggedright
Module
\end{minipage} & \begin{minipage}[b]{\linewidth}\raggedright
Theme
\end{minipage} & \begin{minipage}[b]{\linewidth}\raggedright
Laboratory
\end{minipage} & \begin{minipage}[b]{\linewidth}\raggedright
Gate
\end{minipage} & \begin{minipage}[b]{\linewidth}\raggedleft
Hours
\end{minipage} \\
\midrule\noalign{}
\endhead
\bottomrule\noalign{}
\endlastfoot
01 & M01 & Agent system boundaries and LLM tool semantics & L01 & Lab
review & 10 \\
02 & M01 & Minimal framework-free agent execution & L02 & Lab review &
10 \\
03 & M02 & Framework orchestration: LangGraph and SDK alternatives & L03
& Lab review & 10 \\
04 & M02 & Durable state, checkpoint recovery and human interrupts & L04
& Lab review & 10 \\
05 & M03 & Agent harness control plane, context and environment & L05 &
Lab review & 10 \\
06 & M04 & MCP host/client/server fundamentals and protocol revisions &
L06 & Lab review & 10 \\
07 & M05 & Authenticated GitHub MCP and REST integration & L07 & Lab
review & 10 \\
08 & M06 & Context engineering, retrieval and midterm demonstration &
L08 & Midterm & 10 \\
09 & M07 & Agent threat modeling and adversarial boundaries & L09 & Lab
review & 10 \\
10 & M07 & Policy-as-code, human approval and safe actions & L10 & Lab
review & 10 \\
11 & M08 & Multi-agent delegation and distributed execution & L11 & Lab
review & 10 \\
12 & M09 & Evaluation-driven agent quality engineering & L12 & Lab
review & 10 \\
13 & M10 & Observability, SLOs and operational fault recovery & L13 &
Lab review & 10 \\
14 & M11 & Containerized deployment, CI/CD and workload isolation & L14
& Lab review & 10 \\
15 & M12 & Capstone integration: GitHub AppSec agent (read-only) & L15 &
Capstone & 10 \\
16 & M12 & Capstone hardening, independent review and operational
handoff & L16 & Capstone & 10 \\
\end{longtable}
\end{landscape}

\textbf{Development cadence.} Every weekly lab must be committed in its
own directory with a short README, setup instructions, a failing
negative test, and a run-specific evidence record. Week 4 introduces
durable state; week 6 introduces MCP; week 9 introduces adversarial
prompts; week 10 is the first mock privileged-action workflow. W15--16
remain read-only against any real GitHub repository. Module numbering
follows the 12 requested instructional modules; modules 1, 2, 7 and 12
intentionally span multiple weeks.

\newpage
\subsection{7. Detailed Weekly Syllabus}\label{detailed-weekly-syllabus}

Each weekly plan contains measurable objectives, readings, videos,
concepts, patterns, implementation and testing work, concrete artifacts,
evaluation criteria, and a completion checklist. Required readings are
free unless explicitly labeled optional subscription. Video selections
are curated segments rather than a requirement to finish every
full-length course in a single week.

\subsubsection{Week 01 --- Agent system boundaries and LLM tool
semantics}\label{week-01-agent-system-boundaries-and-llm-tool-semantics}

\textbf{Module:} M01 · \textbf{Prerequisites:} General software
engineering, HTTP/JSON; Python typing primer · \textbf{Hours:} 2
lecture/video + 1.5 reading + 4 coding + 1.5 verification + 1 review =
\textbf{10}.

\textbf{Measurable learning objectives} - Classify responsibilities
among model, framework, harness, MCP, and application governance in six
scenarios. - Implement strict structured response schemas and reject
malformed model outputs. - Construct a threat-boundary diagram
distinguishing proposal from authority.

\textbf{Required reading:} FR01 Unit 1; D06; D07; source PDF sections
1--4.\\
\textbf{Video/course:} YT01; YT05 opening fundamentals.\\
\textbf{Optional depth:} compare corresponding
Microsoft/OpenAI/LangChain reference implementation, or use an
additional LinkedIn Learning lecture if subscribed.

\textbf{Technical study:} Inference versus execution; tool proposal;
structured outputs; JSON Schema; trust boundaries; single versus
multiple agents.\\
\textbf{Engineering pattern / architecture:} Hexagonal architecture and
ports/adapters around a model provider.

\textbf{Laboratory L01 implementation:} Define ModelPort, ToolSpec,
AgentRequest and AgentResult; create a deterministic FakeModel that
returns predefined actions; never execute model-produced source code.\\
\textbf{Test plan:} Property-based or parametrized schema failures;
verify unknown tools are rejected before dispatch.\\
\textbf{Security and reliability emphasis:} Model output is untrusted,
and model reasoning never authorizes an API call.

\textbf{Artifacts:} Boundary diagram, typed schemas, baseline test
report.\\
\textbf{Assessment:} All five layers correctly assigned; three invalid
inputs rejected; no tool invoked without validation.

\textbf{Completion checklist:} {[} {]} reading annotation committed; {[}
{]} demonstration recorded with pinned dependencies; {[} {]} positive
cases passed; {[} {]} negative cases passed; {[} {]} diagram or decision
record committed; {[} {]} reviewer can reproduce the result from one
command.

\subsubsection{Week 02 --- Minimal framework-free agent
execution}\label{week-02-minimal-framework-free-agent-execution}

\textbf{Module:} M01 · \textbf{Prerequisites:} W1; Python async/await
and exception handling · \textbf{Hours:} 2 lecture/video + 1.5 reading +
4 coding + 1.5 verification + 1 review = \textbf{10}.

\textbf{Measurable learning objectives} - Develop a bounded
reason--act--observe state machine without an agent SDK. - Demonstrate
deterministic completion, denial, exhausted-budget, and failure states.
- Prevent duplicate side effects across retries using idempotency keys.

\textbf{Required reading:} D10 (architecture and harness
responsibilities); D07.\\
\textbf{Video/course:} YT02 (00:00--13:18); YT05 harness segments if
verified.\\
\textbf{Optional depth:} compare corresponding
Microsoft/OpenAI/LangChain reference implementation, or use an
additional LinkedIn Learning lecture if subscribed.

\textbf{Technical study:} Finite state machine; asyncio; cancellation;
typed results; retry policy; idempotency; resource budgets.\\
\textbf{Engineering pattern / architecture:} Explicit state machine with
ports/adapters and command dispatcher.

\textbf{Laboratory L02 implementation:} Implement run\_id, max\_steps,
deadline, allowlisted tool catalog, approved input schema and
append-only events.\\
\textbf{Test plan:} Fake timeout, malformed tool argument, unknown tool,
duplicate-call and exact-step-budget fixtures.\\
\textbf{Security and reliability emphasis:} Only pure read-only tools
enabled; all write attempts denied.

\textbf{Artifacts:} Minimal harness package, state-transition table, 10
unit tests.\\
\textbf{Assessment:} No infinite loop, explicit reason for every
terminal state, duplicate invocation prevented.

\textbf{Completion checklist:} {[} {]} reading annotation committed; {[}
{]} demonstration recorded with pinned dependencies; {[} {]} positive
cases passed; {[} {]} negative cases passed; {[} {]} diagram or decision
record committed; {[} {]} reviewer can reproduce the result from one
command.

\subsubsection{Week 03 --- Framework orchestration: LangGraph and SDK
alternatives}\label{week-03-framework-orchestration-langgraph-and-sdk-alternatives}

\textbf{Module:} M02 · \textbf{Prerequisites:} W2; state-machine
understanding · \textbf{Hours:} 2 lecture/video + 1.5 reading + 4 coding
+ 1.5 verification + 1 review = \textbf{10}.

\textbf{Measurable learning objectives} - Implement equivalent routing
using a graph and the handcrafted runner. - Compare LangGraph, OpenAI
Agents SDK and Microsoft Agent Framework against a requirements matrix.
- Demonstrate deterministic graph branch selection with isolated tests.

\textbf{Required reading:} FR02 modules 1--2; D05; D06; D10.\\
\textbf{Video/course:} YT03 00:26:58--01:21:10; LL02 optional.\\
\textbf{Optional depth:} compare corresponding
Microsoft/OpenAI/LangChain reference implementation, or use an
additional LinkedIn Learning lecture if subscribed.

\textbf{Technical study:} Nodes/edges; reducers; event flow; runner;
structured tool adapters; state graph versus chain.\\
\textbf{Engineering pattern / architecture:} Workflow graph plus adapter
layer; avoid framework business-logic coupling.

\textbf{Laboratory L03 implementation:} Build
analysis→retrieve→verify→report nodes with a typed state; maintain same
external contracts as W2.\\
\textbf{Test plan:} Graph branch tests and parity test comparing
identical fixture outputs across implementations.\\
\textbf{Security and reliability emphasis:} Enforce authorization in
application-owned adapters not only model instructions.

\textbf{Artifacts:} Graph implementation, framework ADR-001, parity test
results.\\
\textbf{Assessment:} Graph reproducible; trade-off evaluation has ≥6
criteria and ≥3 technologies.

\textbf{Completion checklist:} {[} {]} reading annotation committed; {[}
{]} demonstration recorded with pinned dependencies; {[} {]} positive
cases passed; {[} {]} negative cases passed; {[} {]} diagram or decision
record committed; {[} {]} reviewer can reproduce the result from one
command.

\subsubsection{Week 04 --- Durable state, checkpoint recovery and human
interrupts}\label{week-04-durable-state-checkpoint-recovery-and-human-interrupts}

\textbf{Module:} M02 · \textbf{Prerequisites:} W3; SQL and serialization
basics · \textbf{Hours:} 2 lecture/video + 1.5 reading + 4 coding + 1.5
verification + 1 review = \textbf{10}.

\textbf{Measurable learning objectives} - Persist execution checkpoints
with monotonic state transitions. - Resume interrupted work without
reissuing a non-idempotent command. - Differentiate operator
intervention from automatic workflow continuation.

\textbf{Required reading:} FR02 modules 2--3; D05; D10.\\
\textbf{Video/course:} YT03 01:22:11--02:08:28.\\
\textbf{Optional depth:} compare corresponding
Microsoft/OpenAI/LangChain reference implementation, or use an
additional LinkedIn Learning lecture if subscribed.

\textbf{Technical study:} Checkpointers; SQLite/PostgreSQL; versioned
state; pause/resume; optimistic concurrency; interrupts.\\
\textbf{Engineering pattern / architecture:} Event sourcing-lite,
durable workflow, approval checkpoint.

\textbf{Laboratory L04 implementation:} Persist state before/after each
tool call and link run\_id to trace\_id; implement pause for approval.\\
\textbf{Test plan:} Kill and restart mid-run; restore from persisted
checkpoint; stale approval rejected.\\
\textbf{Security and reliability emphasis:} Approver identity and policy
version recorded; unapproved resume impossible.

\textbf{Artifacts:} Checkpoint schema, recovery demo, approval boundary
diagram.\\
\textbf{Assessment:} Resumed run produces exactly one durable outcome
and preserves prior evidence.

\textbf{Completion checklist:} {[} {]} reading annotation committed; {[}
{]} demonstration recorded with pinned dependencies; {[} {]} positive
cases passed; {[} {]} negative cases passed; {[} {]} diagram or decision
record committed; {[} {]} reviewer can reproduce the result from one
command.

\subsubsection{Week 05 --- Agent harness control plane, context and
environment}\label{week-05-agent-harness-control-plane-context-and-environment}

\textbf{Module:} M03 · \textbf{Prerequisites:} W2--4; container basics ·
\textbf{Hours:} 2 lecture/video + 1.5 reading + 4 coding + 1.5
verification + 1 review = \textbf{10}.

\textbf{Measurable learning objectives} - Separate tool exposure,
invocation, policy checks, and environment isolation. - Apply token,
cost, iteration and wall-clock budgets. - Demonstrate controlled context
compaction without losing audit provenance.

\textbf{Required reading:} FR03 modules 1--3; D10; D08.\\
\textbf{Video/course:} YT02 09:03--end; YT01.\\
\textbf{Optional depth:} compare corresponding
Microsoft/OpenAI/LangChain reference implementation, or use an
additional LinkedIn Learning lecture if subscribed.

\textbf{Technical study:} Harness architecture; tool catalogs; context
compaction; execution sandbox; skills; hook lifecycle; budgets.\\
\textbf{Engineering pattern / architecture:} Policy enforcement point +
bounded interpreter.

\textbf{Laboratory L05 implementation:} Introduce policy gate, context
store, restricted execution adapter and explicit budget accounting to
W2/W4 harness.\\
\textbf{Test plan:} Adversarial tool output; budget expiry; context
overflow; denied egress; disabled tool.\\
\textbf{Security and reliability emphasis:} No shell/network by default;
tool outputs treated as untrusted.

\textbf{Artifacts:} Harness specification, sandbox config,
controlled-run traces.\\
\textbf{Assessment:} Enforcement remains effective if model prompt is
malicious or contradictory.

\textbf{Completion checklist:} {[} {]} reading annotation committed; {[}
{]} demonstration recorded with pinned dependencies; {[} {]} positive
cases passed; {[} {]} negative cases passed; {[} {]} diagram or decision
record committed; {[} {]} reviewer can reproduce the result from one
command.

\subsubsection{Week 06 --- MCP host/client/server fundamentals and
protocol
revisions}\label{week-06-mcp-hostclientserver-fundamentals-and-protocol-revisions}

\textbf{Module:} M04 · \textbf{Prerequisites:} W1--5; JSON-RPC and
transport basics · \textbf{Hours:} 2 lecture/video + 1.5 reading + 4
coding + 1.5 verification + 1 review = \textbf{10}.

\textbf{Measurable learning objectives} - Diagram exact
host/client/server and tool/result trust boundaries. - Create a local
MCP server and client exposing a tool, resource and prompt. -
Distinguish legacy session initialization from the 2026-07-28 stateless
core.

\textbf{Required reading:} FR05 units 1--2; FR06 MCP Architecture;
D01--D03.\\
\textbf{Video/course:} LL03 optional; YT04 MCP segment (check timestamps
before use).\\
\textbf{Optional depth:} compare corresponding
Microsoft/OpenAI/LangChain reference implementation, or use an
additional LinkedIn Learning lecture if subscribed.

\textbf{Technical study:} JSON-RPC; tools/list and tools/call;
resources/prompts; discovery; stdio; protocol metadata; compatibility
matrix.\\
\textbf{Engineering pattern / architecture:} Capability provider /
anti-corruption layer.

\textbf{Laboratory L06 implementation:} Build mock issue-metadata MCP
server with read-only tools; run dual compatibility experiments against
pinned SDK versions.\\
\textbf{Test plan:} Schema, error response, unsupported version,
capability and malicious-response tests.\\
\textbf{Security and reliability emphasis:} Do not expose entire
conversation; deny unexpected resource identifiers.

\textbf{Artifacts:} MCP architecture diagram, server/client, protocol
compatibility log.\\
\textbf{Assessment:} Five MCP protocol cases documented; correct
capability behavior for chosen SDK version.

\textbf{Completion checklist:} {[} {]} reading annotation committed; {[}
{]} demonstration recorded with pinned dependencies; {[} {]} positive
cases passed; {[} {]} negative cases passed; {[} {]} diagram or decision
record committed; {[} {]} reviewer can reproduce the result from one
command.

\subsubsection{Week 07 --- Authenticated GitHub MCP and REST
integration}\label{week-07-authenticated-github-mcp-and-rest-integration}

\textbf{Module:} M05 · \textbf{Prerequisites:} W6; GitHub APIs and OAuth
concepts · \textbf{Hours:} 2 lecture/video + 1.5 reading + 4 coding +
1.5 verification + 1 review = \textbf{10}.

\textbf{Measurable learning objectives} - Design authorization
boundaries for GitHub App installation tokens. - Expose repo and
code-scanning metadata through read-only MCP tools. - Test expired
credentials, repository boundary violations and rate limiting.

\textbf{Required reading:} FR05 deployed unit; D11; D12; D21; D01.\\
\textbf{Video/course:} LL04 optional; FR06 remote server lesson.\\
\textbf{Optional depth:} compare corresponding
Microsoft/OpenAI/LangChain reference implementation, or use an
additional LinkedIn Learning lecture if subscribed.

\textbf{Technical study:} GitHub App authentication; REST code-scanning
endpoints; token scoping; HTTP OAuth; MCP gateway.\\
\textbf{Engineering pattern / architecture:} Facade with least-privilege
credential broker and API gateway.

\textbf{Laboratory L07 implementation:} Create get\_repo\_metadata,
list\_alerts, get\_alert\_evidence contracts; mock GitHub API first;
optional verified test installation.\\
\textbf{Test plan:} 401/403/404/429; tenant and repo mismatch; network
failures; paging; no write route.\\
\textbf{Security and reliability emphasis:} Separate read and write
identities; no broad personal access tokens in committed configs.

\textbf{Artifacts:} GitHub integration, mocked fixtures, authorization
matrix, SDK pinning note.\\
\textbf{Assessment:} 0 successful forbidden writes; paginated alerts and
429 handled without data loss.

\textbf{Completion checklist:} {[} {]} reading annotation committed; {[}
{]} demonstration recorded with pinned dependencies; {[} {]} positive
cases passed; {[} {]} negative cases passed; {[} {]} diagram or decision
record committed; {[} {]} reviewer can reproduce the result from one
command.

\subsubsection{Week 08 --- Context engineering, retrieval and midterm
demonstration}\label{week-08-context-engineering-retrieval-and-midterm-demonstration}

\textbf{Module:} M06 · \textbf{Prerequisites:} W4--7; elementary
information retrieval · \textbf{Hours:} 2 lecture/video + 1.5 reading +
4 coding + 1.5 verification + 1 review = \textbf{10}.

\textbf{Measurable learning objectives} - Compare lexical and semantic
retrieval quality with a labeled corpus. - Implement context
minimization and evidence provenance. - Present a midterm read-only
investigative agent with recovery and typed MCP tools.

\textbf{Required reading:} FR03 module 3; FR01 unit 3; D05.\\
\textbf{Video/course:} YT04 00:09:36--00:50:40; YT03 03:17:43--end.\\
\textbf{Optional depth:} compare corresponding
Microsoft/OpenAI/LangChain reference implementation, or use an
additional LinkedIn Learning lecture if subscribed.

\textbf{Technical study:} RAG; BM25; embeddings; precision@k; chunking;
memory versus evidence; compaction risks.\\
\textbf{Engineering pattern / architecture:} Retrieval pipeline /
context provenance ledger.

\textbf{Laboratory L08 implementation:} Index 50 synthetic alerts and
code excerpts; retrieve exact commit+rule evidence and produce citation
links.\\
\textbf{Test plan:} Gold-labeled query set; poisoned document;
irrelevant chunks; checkpoint resumption.\\
\textbf{Security and reliability emphasis:} Untrusted retrieved text
never changes policy or tool catalog.

\textbf{Artifacts:} Midterm demo, retrieval benchmark, evidence ledger,
design review.\\
\textbf{Assessment:} All cited artifacts resolve to fixture IDs;
precision@3 measured; no unsupported permission change.

\textbf{Completion checklist:} {[} {]} reading annotation committed; {[}
{]} demonstration recorded with pinned dependencies; {[} {]} positive
cases passed; {[} {]} negative cases passed; {[} {]} diagram or decision
record committed; {[} {]} reviewer can reproduce the result from one
command.

\subsubsection{Week 09 --- Agent threat modeling and adversarial
boundaries}\label{week-09-agent-threat-modeling-and-adversarial-boundaries}

\textbf{Module:} M07 · \textbf{Prerequisites:} W5--8; security
fundamentals · \textbf{Hours:} 2 lecture/video + 1.5 reading + 4 coding
+ 1.5 verification + 1 review = \textbf{10}.

\textbf{Measurable learning objectives} - Construct STRIDE-informed
abuse cases against model/tool/MCP boundaries. - Execute indirect
prompt-injection and tool-poisoning tests. - Identify data exfiltration
and excessive-agency paths with mitigation owners.

\textbf{Required reading:} D13--D16; D08; source PDF section 7.\\
\textbf{Video/course:} LL05 optional.\\
\textbf{Optional depth:} compare corresponding
Microsoft/OpenAI/LangChain reference implementation, or use an
additional LinkedIn Learning lecture if subscribed.

\textbf{Technical study:} Prompt injection; tool response trust; threat
trees; model supply chain; OWASP; MITRE ATLAS.\\
\textbf{Engineering pattern / architecture:} Trust-boundary defense in
depth with explicit security assumptions.

\textbf{Laboratory L09 implementation:} Create 12 attacks: forged system
directions in issue body, malicious MCP tool description, fake approval,
tainted retrieved snippet.\\
\textbf{Test plan:} Adversarial fixtures prove policy denies and audit
redaction still holds.\\
\textbf{Security and reliability emphasis:} Assume repo, issues and tool
outputs are attacker-controlled data.

\textbf{Artifacts:} Threat model, abuse-case tests, control mapping,
residual risk register.\\
\textbf{Assessment:} ≥12 attack cases; 100\% of authorization-bypass
fixtures blocked.

\textbf{Completion checklist:} {[} {]} reading annotation committed; {[}
{]} demonstration recorded with pinned dependencies; {[} {]} positive
cases passed; {[} {]} negative cases passed; {[} {]} diagram or decision
record committed; {[} {]} reviewer can reproduce the result from one
command.

\subsubsection{Week 10 --- Policy-as-code, human approval and safe
actions}\label{week-10-policy-as-code-human-approval-and-safe-actions}

\textbf{Module:} M07 · \textbf{Prerequisites:} W9; policy evaluation
basics · \textbf{Hours:} 2 lecture/video + 1.5 reading + 4 coding + 1.5
verification + 1 review = \textbf{10}.

\textbf{Measurable learning objectives} - Separate AI recommendation
from deterministic policy adjudication. - Write allow/deny rules and
test privileged-action approvals. - Demonstrate time-bound,
replay-resistant approval tokens and fail-closed defaults.

\textbf{Required reading:} D08; D16; D11; D19.\\
\textbf{Video/course:} LL05 optional.\\
\textbf{Optional depth:} compare corresponding
Microsoft/OpenAI/LangChain reference implementation, or use an
additional LinkedIn Learning lecture if subscribed.

\textbf{Technical study:} PDP/PEP; OPA/Rego; approvals; identity;
idempotent state transitions; policy version.\\
\textbf{Engineering pattern / architecture:} Policy decision point (PDP)
and enforcement point (PEP).

\textbf{Laboratory L10 implementation:} Create proposed-disposition
schema and OPA policy; stage a mock dismiss-alert request but never
alter a live alert.\\
\textbf{Test plan:} Positive, negative, missing-identity, stale-policy,
expired-approval, revoked-token, replay tests.\\
\textbf{Security and reliability emphasis:} Sensitive write operations
restricted to a separate mocked service behind authenticated approval.

\textbf{Artifacts:} Policy package, 15 decision-table tests, approver
flow, evidence record.\\
\textbf{Assessment:} All negative gate tests pass; no side-effect if
approval absent.

\textbf{Completion checklist:} {[} {]} reading annotation committed; {[}
{]} demonstration recorded with pinned dependencies; {[} {]} positive
cases passed; {[} {]} negative cases passed; {[} {]} diagram or decision
record committed; {[} {]} reviewer can reproduce the result from one
command.

\subsubsection{Week 11 --- Multi-agent delegation and distributed
execution}\label{week-11-multi-agent-delegation-and-distributed-execution}

\textbf{Module:} M08 · \textbf{Prerequisites:} W4--10; asynchronous
systems · \textbf{Hours:} 2 lecture/video + 1.5 reading + 4 coding + 1.5
verification + 1 review = \textbf{10}.

\textbf{Measurable learning objectives} - Implement supervisor,
investigator, critic and verifier roles. - Compare parallel and
sequential work for correctness and cost. - Prevent untrusted
agent-to-agent data from becoming authorization.

\textbf{Required reading:} FR03 module 4; FR02 module 4; D10.\\
\textbf{Video/course:} LL02 optional.\\
\textbf{Optional depth:} compare corresponding
Microsoft/OpenAI/LangChain reference implementation, or use an
additional LinkedIn Learning lecture if subscribed.

\textbf{Technical study:} Subagents; delegated context; concurrent
tasks; queueing; causal lineage; task budget; cancellation.\\
\textbf{Engineering pattern / architecture:} Supervisor/worker with
isolated capabilities; saga-inspired compensation when needed.

\textbf{Laboratory L11 implementation:} Run
planner→investigator→reviewer with read-only tools and immutable
evidence snapshots; verifier owns completion criteria.\\
\textbf{Test plan:} Out-of-order completions, canceled tasks, worker
failure, bounded delegation depth.\\
\textbf{Security and reliability emphasis:} Each worker receives
narrowed credentials and tool set; supervisor cannot create grants.

\textbf{Artifacts:} Multi-agent graph, execution comparison, concurrency
incident report.\\
\textbf{Assessment:} No worker privilege escalation; replay demonstrates
deterministic verification.

\textbf{Completion checklist:} {[} {]} reading annotation committed; {[}
{]} demonstration recorded with pinned dependencies; {[} {]} positive
cases passed; {[} {]} negative cases passed; {[} {]} diagram or decision
record committed; {[} {]} reviewer can reproduce the result from one
command.

\subsubsection{Week 12 --- Evaluation-driven agent quality
engineering}\label{week-12-evaluation-driven-agent-quality-engineering}

\textbf{Module:} M09 · \textbf{Prerequisites:} W1--11; pytest and
statistics basics · \textbf{Hours:} 2 lecture/video + 1.5 reading + 4
coding + 1.5 verification + 1 review = \textbf{10}.

\textbf{Measurable learning objectives} - Build labeled datasets and
distinguish deterministic checks from judge-model ratings. - Quantify
classification precision, recall and unsupported disposition rate. -
Automate regression evaluation as CI quality gates.

\textbf{Required reading:} FR04 modules 1--3; D09; D11.\\
\textbf{Video/course:} YT05 evaluation portion without unverified
timestamp.\\
\textbf{Optional depth:} compare corresponding
Microsoft/OpenAI/LangChain reference implementation, or use an
additional LinkedIn Learning lecture if subscribed.

\textbf{Technical study:} Golden datasets; contract tests;
precision/recall; confidence calibration; flaky output; evaluator
drift.\\
\textbf{Engineering pattern / architecture:} Testing pyramid plus
separate model-behavior evaluations.

\textbf{Laboratory L12 implementation:} Create 100 synthetic alerts: 40
legitimate, 40 benign candidates, 20 indeterminate; evaluate only
recommendations and evidence.\\
\textbf{Test plan:} Negative tests for unsourced claims; set
precision/recall targets; freeze fixtures; rerun after prompt changes.\\
\textbf{Security and reliability emphasis:} Any unsupported dismissal
suggestion or denied operation is a zero-tolerance event in release
gate.

\textbf{Artifacts:} Evaluation dataset, metrics JSON, CI gate script,
quality report.\\
\textbf{Assessment:} Deterministic policy tests green; metrics
reproducible on fixed fixtures.

\textbf{Completion checklist:} {[} {]} reading annotation committed; {[}
{]} demonstration recorded with pinned dependencies; {[} {]} positive
cases passed; {[} {]} negative cases passed; {[} {]} diagram or decision
record committed; {[} {]} reviewer can reproduce the result from one
command.

\subsubsection{Week 13 --- Observability, SLOs and operational fault
recovery}\label{week-13-observability-slos-and-operational-fault-recovery}

\textbf{Module:} M10 · \textbf{Prerequisites:} W12; logging and metrics
basics · \textbf{Hours:} 2 lecture/video + 1.5 reading + 4 coding + 1.5
verification + 1 review = \textbf{10}.

\textbf{Measurable learning objectives} - Correlate model, harness, MCP,
policy and GitHub operation spans. - Specify error budgets and
measurable SLOs. - Recover safely from network, dependency and storage
faults.

\textbf{Required reading:} FR04 observation/production modules; D09;
D17; D10.\\
\textbf{Video/course:} LL06 optional.\\
\textbf{Optional depth:} compare corresponding
Microsoft/OpenAI/LangChain reference implementation, or use an
additional LinkedIn Learning lecture if subscribed.

\textbf{Technical study:} OpenTelemetry spans; RED metrics; trace
context; SLO/SLA; incident taxonomy; PII redaction.\\
\textbf{Engineering pattern / architecture:} End-to-end distributed
tracing plus graceful degradation.

\textbf{Laboratory L13 implementation:} Emit
trace\_id/run\_id/tool\_call\_id; expose metrics for latency, cost,
denial and replay outcomes; define SLOs.\\
\textbf{Test plan:} Inject API timeout and duplicate delivery; assert no
secret leakage in spans and no duplicate state change.\\
\textbf{Security and reliability emphasis:} Redact source and tokens;
keep immutable audit references separate from raw prompts.

\textbf{Artifacts:} Trace exemplars, dashboards spec, incident runbook,
error budget.\\
\textbf{Assessment:} Every sampled run traces to policy decision;
redaction tests pass.

\textbf{Completion checklist:} {[} {]} reading annotation committed; {[}
{]} demonstration recorded with pinned dependencies; {[} {]} positive
cases passed; {[} {]} negative cases passed; {[} {]} diagram or decision
record committed; {[} {]} reviewer can reproduce the result from one
command.

\subsubsection{Week 14 --- Containerized deployment, CI/CD and workload
isolation}\label{week-14-containerized-deployment-cicd-and-workload-isolation}

\textbf{Module:} M11 · \textbf{Prerequisites:} W5--13; Docker, basic
GitHub Actions · \textbf{Hours:} 2 lecture/video + 1.5 reading + 4
coding + 1.5 verification + 1 review = \textbf{10}.

\textbf{Measurable learning objectives} - Package and scan a
containerized harness service. - Implement static checks, tests, policy
gates and provenance-oriented release controls. - Compare serverless and
long-running runtimes for checkpoint, budget and network limits.

\textbf{Required reading:} D18--D21; D17.\\
\textbf{Video/course:} LL06 optional; FR02 module 6 optional.\\
\textbf{Optional depth:} compare corresponding
Microsoft/OpenAI/LangChain reference implementation, or use an
additional LinkedIn Learning lecture if subscribed.

\textbf{Technical study:} Docker; isolated user; GitHub OIDC; SBOM;
artifacts; IaC; AWS reference deployment.\\
\textbf{Engineering pattern / architecture:} Immutable release pipeline
with least-privilege runtime identity.

\textbf{Laboratory L14 implementation:} Build local Docker Compose demo
with mock providers; optional AWS deployment is designed but not
provisioned.\\
\textbf{Test plan:} Workflow lint, tests, dependency scanning, negative
permission test, deploy/rollback smoke.\\
\textbf{Security and reliability emphasis:} No secrets in logs/images;
nonroot image; no automatic production write scope.

\textbf{Artifacts:} Dockerfile, CI workflow, release checklist, IaC
plan-only specification.\\
\textbf{Assessment:} All gates pass in local/dry-run CI; deployment
design documents rollback.

\textbf{Completion checklist:} {[} {]} reading annotation committed; {[}
{]} demonstration recorded with pinned dependencies; {[} {]} positive
cases passed; {[} {]} negative cases passed; {[} {]} diagram or decision
record committed; {[} {]} reviewer can reproduce the result from one
command.

\subsubsection{Week 15 --- Capstone integration: GitHub AppSec agent
(read-only)}\label{week-15-capstone-integration-github-appsec-agent-read-only}

\textbf{Module:} M12 · \textbf{Prerequisites:} W1--14; passing midterm
and security gates · \textbf{Hours:} 2 lecture/video + 1.5 reading + 4
coding + 1.5 verification + 1 review = \textbf{10}.

\textbf{Measurable learning objectives} - Integrate tools, harness, MCP,
retrieval, policy and tracing into one system. - Produce evidence-linked
recommendations for CodeQL alerts. - Demonstrate read-only behavior
against untrusted repository fixtures.

\textbf{Required reading:} D11; D12; D19; source PDF sections 9--10.\\
\textbf{Video/course:} Optional review YT02, LL05.\\
\textbf{Optional depth:} compare corresponding
Microsoft/OpenAI/LangChain reference implementation, or use an
additional LinkedIn Learning lecture if subscribed.

\textbf{Technical study:} Evidence contracts; CodeQL rule/instance;
repository provenance; recommendation schema; reviewer queue.\\
\textbf{Engineering pattern / architecture:} Read-only investigator
service and explicit recommendation workflow.

\textbf{Laboratory L15 implementation:} Implement
intake→fetch→analyze→verify→report; separate optional mocked approval
service.\\
\textbf{Test plan:} Golden case sets, missing evidence escalation,
cross-repo denial, idempotent restart.\\
\textbf{Security and reliability emphasis:} Forbidden write APIs not
registered in production read-only tool catalog.

\textbf{Artifacts:} Capstone release candidate, 15-item evidence bundle,
reviewer summary.\\
\textbf{Assessment:} All recommendations reproducible and source-backed;
zero live writes.

\textbf{Completion checklist:} {[} {]} reading annotation committed; {[}
{]} demonstration recorded with pinned dependencies; {[} {]} positive
cases passed; {[} {]} negative cases passed; {[} {]} diagram or decision
record committed; {[} {]} reviewer can reproduce the result from one
command.

\subsubsection{Week 16 --- Capstone hardening, independent review and
operational
handoff}\label{week-16-capstone-hardening-independent-review-and-operational-handoff}

\textbf{Module:} M12 · \textbf{Prerequisites:} W15; complete evidence
suite · \textbf{Hours:} 2 lecture/video + 1.5 reading + 4 coding + 1.5
verification + 1 review = \textbf{10}.

\textbf{Measurable learning objectives} - Conduct architecture,
adversarial and recovery reviews against acceptance criteria. - Show
release quality gate and operational rollback procedures. - Produce a
defensible production-readiness decision with residual risks.

\textbf{Required reading:} D13; D14; D19; source PDF sections 7--10.\\
\textbf{Video/course:} LL06 optional.\\
\textbf{Optional depth:} compare corresponding
Microsoft/OpenAI/LangChain reference implementation, or use an
additional LinkedIn Learning lecture if subscribed.

\textbf{Technical study:} Review signoff; residual risk; change
governance; evidence traceability; release checks; handoff.\\
\textbf{Engineering pattern / architecture:} Release readiness review
with independent control evidence.

\textbf{Laboratory L16 implementation:} Freeze version manifest and
evidence bundle; run safety replay plus user-acceptance
demonstrations.\\
\textbf{Test plan:} Reexecute all negative gates; simulated 403, prompt
injection, timeout, restart and manual approval.\\
\textbf{Security and reliability emphasis:} No live dismissal; all
state-changing methods remain mocked until separately approved.

\textbf{Artifacts:} Final engineering report, 15 capstone artifacts,
architecture review decision.\\
\textbf{Assessment:} Safety gates all pass; reviewer can reproduce
build, evaluation and evidence from pinned manifest.

\textbf{Completion checklist:} {[} {]} reading annotation committed; {[}
{]} demonstration recorded with pinned dependencies; {[} {]} positive
cases passed; {[} {]} negative cases passed; {[} {]} diagram or decision
record committed; {[} {]} reviewer can reproduce the result from one
command.

\subsection{8. Module-Level Technical Study
Guides}\label{module-level-technical-study-guides}

These are compact topic maps and engineering checkpoints to accompany
the weekly plans; the official vendor references, listed in section 23,
remain the source of API-level truth.

\begin{landscape}
\begin{longtable}[]{@{}
  >{\raggedright\arraybackslash}p{(\columnwidth - 6\tabcolsep) * \real{0.1500}}
  >{\raggedright\arraybackslash}p{(\columnwidth - 6\tabcolsep) * \real{0.4300}}
  >{\raggedright\arraybackslash}p{(\columnwidth - 6\tabcolsep) * \real{0.1900}}
  >{\raggedright\arraybackslash}p{(\columnwidth - 6\tabcolsep) * \real{0.2300}}@{}}
\toprule\noalign{}
\begin{minipage}[b]{\linewidth}\raggedright
Module
\end{minipage} & \begin{minipage}[b]{\linewidth}\raggedright
Technical knowledge to master
\end{minipage} & \begin{minipage}[b]{\linewidth}\raggedright
Design exercise
\end{minipage} & \begin{minipage}[b]{\linewidth}\raggedright
Evidence of mastery
\end{minipage} \\
\midrule\noalign{}
\endhead
\bottomrule\noalign{}
\endlastfoot
M01 -- Agent foundations & Structured responses, roles, tool request
parsing, loop exit; distinguish inference from execution & ModelPort and
allowlisted ToolSpec & Malformed proposals never cause dispatch \\
M02 -- Framework engineering & State graphs, routing, node reducers,
handoffs, checkpoint compatibility & ADR: framework and state model
selection & Equivalent fixture behavior with/without framework \\
M03 -- Harness runtime & Context budgets, cancellation, event logs,
approval interruption, sandbox boundaries & Runtime state-transition
specification & Deterministic \texttt{blocked} and
\texttt{budget\_exhausted} traces \\
M04 -- MCP fundamentals & Host/client/server, JSON-RPC,
tool/resource/prompt semantics, transport, revision negotiation &
Compare session-based client and 2026-07-28 behavior & Versioned
protocol test report \\
M05 -- MCP integration & HTTP authorization, GitHub Apps, schema and API
pagination, remote gateway & Define read-only \texttt{list\_alerts}
capability contract & 401/403/429 tests and repo boundary rejection \\
M06 -- Context/memory & RAG, chunks, BM25, embeddings, context
compaction, persistent evidence & Evidence retrieval policy &
Precision@k and provenance for each claim \\
M07 -- Security/guardrails & Prompt injection, tool poisoning, OWASP,
NIST, OPA, approvals and logging & Abuse-case inventory and PEP/PDP
split & No privileged effect under adversarial conditions \\
M08 -- Multiagent & Supervisor/worker, typed task handoffs, parallel
cancellation and isolated state & Bounded delegation topology & Worker
failure leaves recoverable complete audit \\
M09 -- Evaluation & Golden fixtures, regression, precision/recall, judge
bias, deterministic policies & Evaluator design record & Reproducible
quality delta across two releases \\
M10 -- Reliability & OpenTelemetry, trace propagation, metrics, SLOs,
alerts and incident review & Operational telemetry schema & Drill-down
from recommendation to tool/policy spans \\
M11 -- Deployment & Containers, CI/CD, SBOM, secrets, OIDC, AWS decision
trade-offs & Deployment ADR + rollback plan & CI blocks failing security
gates \\
M12 -- AppSec capstone & CodeQL alert evidence, authorized repository
boundaries, reviewer recommendations & End-to-end investigator
architecture & Versioned review packet and independent re-execution \\
\end{longtable}
\end{landscape}

\subsubsection{MCP compatibility lab: old versus July
2026}\label{mcp-compatibility-lab-old-versus-july-2026}

The published 2026-07-28 MCP specification uses a stateless protocol
core, self-contained requests, optional \texttt{server/discover}, and
revised multi-round-trip interactions. Older tutorials may assume an
\texttt{initialize} handshake and session identifier. \textbf{Do not mix
wire-level examples across versions.} Each implementation records the
selected SDK package version, spec revision, transport, supported
methods, discovery behavior and authorization setup. Maintain one test
fixture for the supported modern version and one explicit
\texttt{unsupported\_version} fixture. The 2026 Java modern-server API
is labeled preview; use its documented API and do not claim broad
production compatibility without tests. Refer to D01--D04.

\subsubsection{Cross-language assignment
options}\label{cross-language-assignment-options}

\begin{itemize}
\tightlist
\item
  \textbf{Python core:} typed protocols, \texttt{asyncio},
  \texttt{pytest}, SQLite, Python MCP SDK; completed in the main course.
\item
  \textbf{TypeScript extension:} reuse JSON Schemas, implement a Node
  MCP gateway and contract tests; work from TypeScript SDK v2 migration
  notes rather than auto-translating legacy v1 code.
\item
  \textbf{Java extension:} implement a Spring Boot client façade, signed
  token verification, JUnit contract suite and secure exception mapping;
  prefer stable SDK surfaces and explicitly label Java modern-server
  preview use.
\item
  \textbf{Deployment extension:} AWS infrastructure is architecture-only
  unless the instructor/operator separately authorizes resources,
  budgets, secrets and cleanup.
\end{itemize}

\subsection{9. Verified LinkedIn Learning Course
Catalog}\label{verified-linkedin-learning-course-catalog}

LinkedIn Learning access is \textbf{subscription-based}, with a
one-month trial advertised on the verified course pages. A certificate
for a LinkedIn course or path documents course completion; it does not
establish competence on this curriculum's independent engineering safety
gates. Some listed tutorials predate the July 2026 MCP change.

\subsubsection{LL01: Building Agentic AI Systems for
Developers}\label{ll01-building-agentic-ai-systems-for-developers}

\begin{itemize}
\tightlist
\item
  \textbf{Provider/instructor:} LinkedIn Learning; \textbf{access:}
  Subscription / free trial; \textbf{listed length/date:} 11 courses; 16
  h; path updated 2026-05-01.
\item
  \textbf{Official listing:}
  \href{https://www.linkedin.com/learning/paths/building-agentic-ai-systems-for-developers}{Building
  Agentic AI Systems for Developers}.
\item
  \textbf{Relevant coverage:} Developer path: tools, data, MCP,
  evaluation, safety and operations. \textbf{Curriculum placement:}
  Optional umbrella track; contents subject to announced 2026-11-16
  update.
\item
  \textbf{Prerequisites/level:} baseline Python/API knowledge is
  recommended for technical courses; provider labels LL02/LL05/LL06
  intermediate, LL03 beginner, and LL04 general. The 11-course path
  spans beginner/intermediate material.
\item
  \textbf{Applied assignment:} annotate one agent boundary or implement
  an equivalent local mock; never accept course-demo authorization as
  sufficient for production.
\item
  \textbf{Limitations / validation:} Verified provider page; curriculum
  may change. Always check current access status before enrollment.
\end{itemize}

\subsubsection{LL02: Building AI Agents: From Zero to Production by
Microsoft}\label{ll02-building-ai-agents-from-zero-to-production-by-microsoft}

\begin{itemize}
\tightlist
\item
  \textbf{Provider/instructor:} Microsoft Developer; \textbf{access:}
  Subscription / free trial; \textbf{listed length/date:} 50 min;
  released 2025-12-11.
\item
  \textbf{Official listing:}
  \href{https://www.linkedin.com/learning/building-ai-agents-from-zero-to-production-by-microsoft}{Building
  AI Agents: From Zero to Production by Microsoft}.
\item
  \textbf{Relevant coverage:} Microsoft Agent Framework; multiagent
  development-onboarding application. \textbf{Curriculum placement:}
  W3--4, W11.
\item
  \textbf{Prerequisites/level:} baseline Python/API knowledge is
  recommended for technical courses; provider labels LL02/LL05/LL06
  intermediate, LL03 beginner, and LL04 general. The 11-course path
  spans beginner/intermediate material.
\item
  \textbf{Applied assignment:} annotate one agent boundary or implement
  an equivalent local mock; never accept course-demo authorization as
  sufficient for production.
\item
  \textbf{Limitations / validation:} Verified provider page. Always
  check current access status before enrollment.
\end{itemize}

\subsubsection{LL03: Introduction to Model Context Protocol by
Anthropic}\label{ll03-introduction-to-model-context-protocol-by-anthropic}

\begin{itemize}
\tightlist
\item
  \textbf{Provider/instructor:} Anthropic; \textbf{access:} Subscription
  / free trial; \textbf{listed length/date:} 61 min; released
  2026-03-30.
\item
  \textbf{Official listing:}
  \href{https://www.linkedin.com/learning/introduction-to-model-context-protocol-by-anthropic}{Introduction
  to Model Context Protocol by Anthropic}.
\item
  \textbf{Relevant coverage:} Python MCP server/client, tools,
  resources, prompts, Inspector. \textbf{Curriculum placement:} W6.
\item
  \textbf{Prerequisites/level:} baseline Python/API knowledge is
  recommended for technical courses; provider labels LL02/LL05/LL06
  intermediate, LL03 beginner, and LL04 general. The 11-course path
  spans beginner/intermediate material.
\item
  \textbf{Applied assignment:} annotate one agent boundary or implement
  an equivalent local mock; never accept course-demo authorization as
  sufficient for production.
\item
  \textbf{Limitations / validation:} Verified provider page; legacy
  protocol examples must be compared to 2026-07-28. Always check current
  access status before enrollment.
\end{itemize}

\subsubsection{LL04: Unboxing AI: Build a Remote MCP Server from Zero to
Deployed with
OAuth}\label{ll04-unboxing-ai-build-a-remote-mcp-server-from-zero-to-deployed-with-oauth}

\begin{itemize}
\tightlist
\item
  \textbf{Provider/instructor:} Morten Rand-Hendriksen; \textbf{access:}
  Subscription / free trial; \textbf{listed length/date:} 2 h 55 min;
  released 2025-07-31.
\item
  \textbf{Official listing:}
  \href{https://www.linkedin.com/learning/unboxing-ai-build-a-remote-mcp-server-from-zero-to-deployed-with-oauth}{Unboxing
  AI: Build a Remote MCP Server from Zero to Deployed with OAuth}.
\item
  \textbf{Relevant coverage:} Remote MCP deployment, authorization,
  OAuth, elicitation. \textbf{Curriculum placement:} W7.
\item
  \textbf{Prerequisites/level:} baseline Python/API knowledge is
  recommended for technical courses; provider labels LL02/LL05/LL06
  intermediate, LL03 beginner, and LL04 general. The 11-course path
  spans beginner/intermediate material.
\item
  \textbf{Applied assignment:} annotate one agent boundary or implement
  an equivalent local mock; never accept course-demo authorization as
  sufficient for production.
\item
  \textbf{Limitations / validation:} Verified provider page; older MCP
  semantics likely. Always check current access status before
  enrollment.
\end{itemize}

\subsubsection{LL05: Governing AI Agents: Visibility and
Control}\label{ll05-governing-ai-agents-visibility-and-control}

\begin{itemize}
\tightlist
\item
  \textbf{Provider/instructor:} Kesha Williams; \textbf{access:}
  Subscription / free trial; \textbf{listed length/date:} 53 min;
  released 2026-01-27.
\item
  \textbf{Official listing:}
  \href{https://www.linkedin.com/learning/governing-ai-agents-visibility-and-control}{Governing
  AI Agents: Visibility and Control}.
\item
  \textbf{Relevant coverage:} Python guardrails, inventories, action
  logs, approval. \textbf{Curriculum placement:} W9--10.
\item
  \textbf{Prerequisites/level:} baseline Python/API knowledge is
  recommended for technical courses; provider labels LL02/LL05/LL06
  intermediate, LL03 beginner, and LL04 general. The 11-course path
  spans beginner/intermediate material.
\item
  \textbf{Applied assignment:} annotate one agent boundary or implement
  an equivalent local mock; never accept course-demo authorization as
  sufficient for production.
\item
  \textbf{Limitations / validation:} Verified provider page. Always
  check current access status before enrollment.
\end{itemize}

\subsubsection{LL06: Operating AI Agents: Failure and
Recovery}\label{ll06-operating-ai-agents-failure-and-recovery}

\begin{itemize}
\tightlist
\item
  \textbf{Provider/instructor:} Kesha Williams; \textbf{access:}
  Subscription / free trial; \textbf{listed length/date:} 41 min;
  released 2026-02-04.
\item
  \textbf{Official listing:}
  \href{https://www.linkedin.com/learning/operating-ai-agents-failure-and-recovery}{Operating
  AI Agents: Failure and Recovery}.
\item
  \textbf{Relevant coverage:} Failure detection, assessment, recovery,
  post-recovery verification. \textbf{Curriculum placement:} W13--14.
\item
  \textbf{Prerequisites/level:} baseline Python/API knowledge is
  recommended for technical courses; provider labels LL02/LL05/LL06
  intermediate, LL03 beginner, and LL04 general. The 11-course path
  spans beginner/intermediate material.
\item
  \textbf{Applied assignment:} annotate one agent boundary or implement
  an equivalent local mock; never accept course-demo authorization as
  sufficient for production.
\item
  \textbf{Limitations / validation:} Verified provider page. Always
  check current access status before enrollment.
\end{itemize}

\textbf{Selection strategy.} LL01 is a comprehensive subscription
pathway but exceeds the course's 36-hour lecture budget if consumed in
full on top of everything else. Use it as an \textbf{alternative route}:
replace, rather than duplicate, selected free lectures. LL03 and LL04
are conceptual/implementation companions to FR05 and FR06, not
replacements for current MCP compatibility tests. LL05/LL06 are valuable
because most agent tutorials underemphasize operations.

\subsection{10. Verified Free Course
Catalog}\label{verified-free-course-catalog}

The sources below were inspected on their official provider pages.
``Free'' applies to instructional materials, not automatically to model
API calls, cloud deployment, hosted observability, or certificate/graded
features.

\subsubsection{FR01: AI Agents Course}\label{fr01-ai-agents-course}

\begin{itemize}
\tightlist
\item
  \textbf{Publisher/instructor:} Hugging Face. \textbf{Access:} Free.
  \textbf{Length:} Self-paced; units 0--4 plus bonuses; duration not
  published here.
\item
  \textbf{Course:}
  \href{https://huggingface.co/learn/agents-course/unit0/introduction}{AI
  Agents Course}. \textbf{Primary weeks:} W1--3.
\item
  \textbf{Learning goal:} Agent fundamentals, frameworks, use cases and
  evaluation.
\item
  \textbf{Verified syllabus alignment and practical notes:} Verified
  syllabus units: agent fundamentals, frameworks, use cases, final
  assignment, bonus observability/evaluation. Languages/frameworks:
  Python, smolagents, LangGraph, LlamaIndex. Prerequisite: basic
  programming, introductory LLMs. Exercise: adapt a simple agent and
  test its tools.
\item
  \textbf{Caveat:} Verified official course.
\end{itemize}

\subsubsection{FR02: Introduction to LangGraph --
Python}\label{fr02-introduction-to-langgraph-python}

\begin{itemize}
\tightlist
\item
  \textbf{Publisher/instructor:} LangChain Academy. \textbf{Access:}
  Free. \textbf{Length:} 55 lessons; 6 h video.
\item
  \textbf{Course:}
  \href{https://academy.langchain.com/courses/intro-to-langgraph}{Introduction
  to LangGraph -- Python}. \textbf{Primary weeks:} W3--4.
\item
  \textbf{Learning goal:} Graph nodes, routing, state, reducers,
  checkpointing, human feedback, deployment.
\item
  \textbf{Verified syllabus alignment and practical notes:} Verified
  chapters: M1 Simple Graph, Router, Agent; M2 State Schema, Reducers,
  Trim and Summarize; M3 Breakpoints, Time Travel; M4
  Parallelization/Subgraphs; M5 Long-Term Memory; M6 Deployment. Python
  primary. Prerequisite: programming and Python. Exercise: persist a
  graph and replay an interrupted fixture.
\item
  \textbf{Caveat:} Verified official course.
\end{itemize}

\subsubsection{FR03: Introduction to Deep
Agents}\label{fr03-introduction-to-deep-agents}

\begin{itemize}
\tightlist
\item
  \textbf{Publisher/instructor:} LangChain Academy. \textbf{Access:}
  Free. \textbf{Length:} 46 lessons; 0.5 h video as listed.
\item
  \textbf{Course:}
  \href{https://academy.langchain.com/courses/foundation-introduction-to-deepagents}{Introduction
  to Deep Agents}. \textbf{Primary weeks:} W5, W8, W11.
\item
  \textbf{Learning goal:} Harness-oriented execution environment,
  filesystems, sandbox, MCP, memory, subagents.
\item
  \textbf{Verified syllabus alignment and practical notes:} Verified
  lessons: Tools, MCP, Messages/Threads/Checkpointers, HITL, Filesystem
  Backends, Sandboxes and LocalShell, Summarization/Offload, Skills,
  Memory, Delegation. Python/TypeScript setup. Prerequisite: framework
  basics. Exercise: constrain a deep agent to isolated read-only
  context.
\item
  \textbf{Caveat:} Verified official course; lessons include text.
\end{itemize}

\subsubsection{FR04: Building Reliable
Agents}\label{fr04-building-reliable-agents}

\begin{itemize}
\tightlist
\item
  \textbf{Publisher/instructor:} LangChain Academy. \textbf{Access:}
  Free. \textbf{Length:} 29 lessons; video duration listed as 0 h.
\item
  \textbf{Course:}
  \href{https://academy.langchain.com/courses/building-reliable-agents}{Building
  Reliable Agents}. \textbf{Primary weeks:} W12--13.
\item
  \textbf{Learning goal:} Observability, tracing, datasets,
  code/LLM-based evaluations, production.
\item
  \textbf{Verified syllabus alignment and practical notes:} Verified
  modules: Observation (LangSmith tracing), Evaluation (datasets,
  code-based evaluator, LLM-as-judge, pairwise), Moving Towards
  Production (online evals). Python/TypeScript setup. Prerequisite: a
  working agent. Exercise: score frozen fixtures. LangSmith hosted
  features may require an account and have separate plan restrictions.
\item
  \textbf{Caveat:} Verified official course; not a 29-video program.
\end{itemize}

\subsubsection{FR05: Model Context Protocol (MCP)
Course}\label{fr05-model-context-protocol-mcp-course}

\begin{itemize}
\tightlist
\item
  \textbf{Publisher/instructor:} Hugging Face and Anthropic.
  \textbf{Access:} Free. \textbf{Length:} Unit-based, recommended 3--4
  h/week for each unit.
\item
  \textbf{Course:}
  \href{https://huggingface.co/learn/mcp-course/unit0/introduction}{Model
  Context Protocol (MCP) Course}. \textbf{Primary weeks:} W6--7.
\item
  \textbf{Learning goal:} MCP host-client-server, SDK usage, deployed
  projects, certificates.
\item
  \textbf{Verified syllabus alignment and practical notes:} Verified
  units: 1 Fundamentals, 2 End-to-end MCP use case, 3 Deployed use case,
  4 Bonus. Python or TypeScript; baseline APIs and LLMs. Exercise:
  implement a server with tool/resource/prompt and verify with
  Inspector; revisit on 2026-07-28 compatibility.
\item
  \textbf{Caveat:} Verified official course; check SDK version.
\end{itemize}

\subsubsection{FR06: MCP: Build Rich-Context AI Apps with
Anthropic}\label{fr06-mcp-build-rich-context-ai-apps-with-anthropic}

\begin{itemize}
\tightlist
\item
  \textbf{Publisher/instructor:} Elie Schoppik / Anthropic.
  \textbf{Access:} Free for limited beta; some PRO features.
  \textbf{Length:} 1 h 38 min; 11 video lessons; 6 code examples.
\item
  \textbf{Course:}
  \href{https://www.deeplearning.ai/courses/mcp-build-rich-context-ai-apps-with-anthropic}{MCP:
  Build Rich-Context AI Apps with Anthropic}. \textbf{Primary weeks:}
  W6--7.
\item
  \textbf{Learning goal:} MCP tools/resources/prompts, local clients,
  remote server.
\item
  \textbf{Verified syllabus alignment and practical notes:} Verified
  lessons: Why MCP, Architecture, Chatbot Example, Creating Server,
  Creating Client, Connecting Reference Servers, Prompt/Resource
  Features, Remote Deployment. Python examples. Prerequisite: Python and
  basic LLM workflows. Some graded/pro credential features are not free.
\item
  \textbf{Caveat:} Verified course page; free access may change.
\end{itemize}

\subsection{11. Verified YouTube Video and Tutorial
Catalog}\label{verified-youtube-video-and-tutorial-catalog}

YouTube materials supplement, rather than substitute for, current
official SDK and security documentation. ``Verified'' below means the
video title/link or publisher listing was inspected directly;
``corroborated'' means a creator repository or external chapter index
supports the URL; ``provisional'' means sufficient source evidence was
not available for full metadata validation.

\subsubsection{YT01: What is an agent harness? Explained
fast}\label{yt01-what-is-an-agent-harness-explained-fast}

\begin{itemize}
\tightlist
\item
  \textbf{Publisher:} LangChain · \textbf{access:} Free ·
  \textbf{duration/date:} Released 2026-09-11; short explainer.
\item
  \textbf{Video URL:} \url{https://www.youtube.com/watch?v=7dFvU9zuv0k}.
  \textbf{Relevant weeks:} W1, W5.
\item
  \textbf{Topics:} Filesystem, subagents, skills and harness purpose.
  \textbf{Verification:} Verified video listing.
\item
  \textbf{Chapters / implementation exercise:} No chapter markers
  needed: short conceptual video. Implement an independent validation or
  test instead of copying code without review.
\end{itemize}

\subsubsection{YT02: What is an Agent Harness? and How to build a great
one!}\label{yt02-what-is-an-agent-harness-and-how-to-build-a-great-one}

\begin{itemize}
\tightlist
\item
  \textbf{Publisher:} Prompt Engineering · \textbf{access:} Free ·
  \textbf{duration/date:} Released 2026-04-30; video chapters 00:00,
  01:50, 03:19, 13:04.
\item
  \textbf{Video URL:} \url{https://www.youtube.com/watch?v=nWzXyjXCoCE}.
  \textbf{Relevant weeks:} W2, W5.
\item
  \textbf{Topics:} Nine harness components and minimal Python
  implementation. \textbf{Verification:} Verified video and chapter
  markers.
\item
  \textbf{Chapters / implementation exercise:} Verified video chapters:
  00:00 What Is a Harness; 01:50 Harness vs Frameworks; 03:19 Nine
  Components; 09:03 Memory, Prompts and Hooks; 11:41 Permissions and
  Safety; 13:04 Build a Mini Harness. Implement an independent
  validation or test instead of copying code without review.
\end{itemize}

\subsubsection{YT03: LangGraph Complete Course for Beginners -- Complex
AI Agents with
Python}\label{yt03-langgraph-complete-course-for-beginners-complex-ai-agents-with-python}

\begin{itemize}
\tightlist
\item
  \textbf{Publisher:} freeCodeCamp / Vaibhav Mehra · \textbf{access:}
  Free · \textbf{duration/date:} \textasciitilde3 h 35 min; 2025-05-20
  companion article.
\item
  \textbf{Video URL:} \url{https://www.youtube.com/watch?v=jGg_1h0qzaM}.
  \textbf{Relevant weeks:} W3, W8.
\item
  \textbf{Topics:} State graphs, agent exercises, RAG agent and tests.
  \textbf{Verification:} Direct video corroborated by author code
  repository; timestamps available.
\item
  \textbf{Chapters / implementation exercise:} Companion index: 00:26:58
  Type Annotations; 00:34:09 Elements; 00:46:42 Agent 1 code; 01:56:35
  AI Agent; 03:17:43 RAG Agent. These timestamps derive from a published
  companion article, not a freshly retrieved video transcript. Implement
  an independent validation or test instead of copying code without
  review.
\end{itemize}

\subsubsection{YT04: RAG \& MCP Fundamentals -- A Hands-On Crash
Course}\label{yt04-rag-mcp-fundamentals-a-hands-on-crash-course}

\begin{itemize}
\tightlist
\item
  \textbf{Publisher:} freeCodeCamp · \textbf{access:} Free ·
  \textbf{duration/date:} Published companion 2026-01-22; total duration
  unverified.
\item
  \textbf{Video URL:} \url{https://www.youtube.com/watch?v=I7_WXKhyGms}.
  \textbf{Relevant weeks:} W6, W8.
\item
  \textbf{Topics:} BM25, embeddings, chunking, semantic retrieval, MCP.
  \textbf{Verification:} URL and chapters corroborated by secondary
  article; recheck in YouTube.
\item
  \textbf{Chapters / implementation exercise:} Companion index: 00:09:36
  lexical search; 00:27:08 semantic search; 00:41:40 chunking; 00:50:40
  full RAG pipeline. Verify later MCP timestamps directly in player.
  Implement an independent validation or test instead of copying code
  without review.
\end{itemize}

\subsubsection{YT05: Complete Agentic AI Course In 10 Hours ---
LangChain, LangGraph, RAG, Guardrails,
Evals}\label{yt05-complete-agentic-ai-course-in-10-hours-langchain-langgraph-rag-guardrails-evals}

\begin{itemize}
\tightlist
\item
  \textbf{Publisher:} Krish Naik · \textbf{access:} Free ·
  \textbf{duration/date:} Title says 10 hours; exact runtime not
  independently confirmed.
\item
  \textbf{Video URL:} \url{https://www.youtube.com/watch?v=rV3HJ4LEZ7k}.
  \textbf{Relevant weeks:} W2--3, W9, W12.
\item
  \textbf{Topics:} LangGraph, Deep Agents, RAG, guardrails and
  evaluations. \textbf{Verification:} Video title verified; do not infer
  precise timestamps.
\item
  \textbf{Chapters / implementation exercise:} No reliable chapter
  markers were retrieved; use by topic, not quoted timestamps. Implement
  an independent validation or test instead of copying code without
  review.
\end{itemize}

\subsubsection{YT06: Agentic AI Complete Course for
Beginners}\label{yt06-agentic-ai-complete-course-for-beginners}

\begin{itemize}
\tightlist
\item
  \textbf{Publisher:} freeCodeCamp / Bappy · \textbf{access:} Free ·
  \textbf{duration/date:} Length not independently verified.
\item
  \textbf{Video URL:} \url{https://www.youtube.com/watch?v=Zy7EXDONlTY}.
  \textbf{Relevant weeks:} Optional W1--4.
\item
  \textbf{Topics:} Async Python, Pydantic, LangChain, LangGraph,
  multiagent. \textbf{Verification:} URL and chapter listing located on
  secondary mirror; metadata provisional.
\item
  \textbf{Chapters / implementation exercise:} A third-party mirror
  supplied chapter metadata; use this as an optional discovery link, not
  a required graded course. Implement an independent validation or test
  instead of copying code without review.
\end{itemize}

\subsection{12. Course-to-Learning-Objective Traceability
Matrix}\label{course-to-learning-objective-traceability-matrix}

\begin{landscape}
\begin{longtable}[]{@{}
  >{\raggedright\arraybackslash}p{(\columnwidth - 8\tabcolsep) * \real{0.0900}}
  >{\raggedright\arraybackslash}p{(\columnwidth - 8\tabcolsep) * \real{0.1100}}
  >{\raggedright\arraybackslash}p{(\columnwidth - 8\tabcolsep) * \real{0.2000}}
  >{\raggedright\arraybackslash}p{(\columnwidth - 8\tabcolsep) * \real{0.2700}}
  >{\raggedright\arraybackslash}p{(\columnwidth - 8\tabcolsep) * \real{0.3300}}@{}}
\toprule\noalign{}
\begin{minipage}[b]{\linewidth}\raggedright
Outcome
\end{minipage} & \begin{minipage}[b]{\linewidth}\raggedright
Instructional weeks
\end{minipage} & \begin{minipage}[b]{\linewidth}\raggedright
Required training
\end{minipage} & \begin{minipage}[b]{\linewidth}\raggedright
Optional subscription or YouTube
\end{minipage} & \begin{minipage}[b]{\linewidth}\raggedright
Primary evaluated evidence
\end{minipage} \\
\midrule\noalign{}
\endhead
\bottomrule\noalign{}
\endlastfoot
LO01 & 1 & FR01; D06--D07 & YT01, LL01 & \texttt{boundaries.dot}, schema
tests \\
LO02 & 2 & D07, D10 & YT02, YT05 & Bounded runner and traces \\
LO03 & 3--4 & FR02, D05 & LL02, YT03 & Graph parity/restore suite \\
LO04 & 4--5 & FR03, D10 & YT01, YT02 & Harness budget/compaction
tests \\
LO05 & 6--7 & FR05, FR06; D01--D04 & LL03, LL04 & SDK compatibility and
GitHub contracts \\
LO06 & 8 & FR01 unit 3; FR03 & YT04 & Retrieval metrics and
provenance \\
LO07 & 9 & D13--D15 & LL05 & Threat model and injection tests \\
LO08 & 10 & D08, D16 & LL05 & OPA decision table and approval tests \\
LO09 & 11 & FR03 module 4; FR02 module 4 & LL02 & Isolated worker
orchestration \\
LO10 & 12 & FR04; D09 & YT05 & Evaluation dataset and CI report \\
LO11 & 13 & FR04; D17 & LL06 & Telemetry/runbook/fault tests \\
LO12 & 14 & D18--D21 & LL06 & Container release and CI gates \\
LO13 & 15 & D11, D12 & LL05 & Evidence-backed recommendation packet \\
LO14 & 16 & D13--D19 & LL06 & Signed/dated review decision \\
\end{longtable}
\end{landscape}

\textbf{Course resource substitution rule:} paid content is never
essential. For each LinkedIn Learning item the listed free official
documentation and FR/YT material provides an alternate path to the same
tested competence. Official product docs are used for API correctness;
video examples are treated as dated pedagogical snapshots.

\subsection{13. Programming Laboratory
Catalog}\label{programming-laboratory-catalog}

All 16 laboratories are substantial programming exercises, not viewing
assignments. Their shared \textbf{submission contract} is: source and
pinned dependency manifest; \texttt{README.md}; diagram; typed
interfaces; automated positive/negative tests; security note;
verification log; and \texttt{evidence.json} with run ID, git SHA,
timestamp, tool/SDK versions, tester and result. Each laboratory begins
with a \texttt{FakeModel} and fake external service and must remain
runnable offline.

\textbf{Standard repository skeleton} (extend as labs progress):

\begin{Shaded}
\begin{Highlighting}[]
\NormalTok{agent{-}engineering{-}curriculum/}
\NormalTok{  pyproject.toml                  \# pinned Python dependencies}
\NormalTok{  Makefile                        \# setup, test, lint, verify}
\NormalTok{  src/agent\_course/\{contracts,harness,framework,mcp,policy,retrieval,telemetry\}/}
\NormalTok{  labs/L01/ ... labs/L16/         \# isolated exercise entry points and README files}
\NormalTok{  fixtures/\{model,github,alerts,attack\_cases\}/}
\NormalTok{  tests/\{unit,contract,integration,adversarial,evaluation\}/}
\NormalTok{  policy/                          \# OPA/Rego + decision tests}
\NormalTok{  docs/\{architecture,adrs,threat{-}models,runbooks,evidence\}/}
\NormalTok{  .github/workflows/agent{-}course{-}ci.yml}
\end{Highlighting}
\end{Shaded}

\textbf{One-time setup} (illustrative; learners pin actual compatible
versions during W1 and commit the lockfile):

\begin{Shaded}
\begin{Highlighting}[]
\ExtensionTok{python3} \AttributeTok{{-}m}\NormalTok{ venv .venv}
\BuiltInTok{.}\NormalTok{ .venv/bin/activate}
\ExtensionTok{python} \AttributeTok{{-}m}\NormalTok{ pip install }\AttributeTok{{-}U}\NormalTok{ pip}
\ExtensionTok{python} \AttributeTok{{-}m}\NormalTok{ pip install }\AttributeTok{{-}e} \StringTok{\textquotesingle{}.[dev]\textquotesingle{}}
\ExtensionTok{python} \AttributeTok{{-}m}\NormalTok{ pytest }\AttributeTok{{-}q}
\ExtensionTok{python} \AttributeTok{{-}m}\NormalTok{ compileall }\AttributeTok{{-}q}\NormalTok{ src}
\end{Highlighting}
\end{Shaded}

\textbf{Engineering policy:} do not allow an unresolved dependency to be
auto-upgraded silently. Record \texttt{python\ -\/-version}, installed
SDK versions, MCP revision and GitHub REST API version in a
reproducibility manifest. An optional local mock inference adapter is
sufficient; real API calls are graded only after mock tests pass.

\subsubsection{L01 --- Agent system boundaries and LLM tool
semantics}\label{l01-agent-system-boundaries-and-llm-tool-semantics}

\textbf{Purpose and prerequisites:} Execute the Week 01 engineering
outcome with the prerequisites stated in Week 01. \textbf{Target time:}
4 h implementation plus 1.5 h verification (other course work handled in
the weekly allocation).

\textbf{Learning objective:} Classify responsibilities among model,
framework, harness, MCP, and application governance in six scenarios.
\textbf{Architecture diagram (logical):}

\begin{Shaded}
\begin{Highlighting}[]
\NormalTok{User request → schema validator → FakeModel → proposal}
\NormalTok{validator → report}
\end{Highlighting}
\end{Shaded}

\textbf{Required tools:} Python 3, pytest, typed model stub, fixtures
and Git; add only the week-specific dependencies (MCP SDK in L06--07,
OPA in L10, OpenTelemetry in L13, Docker in L14); pin each one in the
lockfile.

\textbf{Starter files / repository structure:} create a directory
\texttt{labs/L01/} containing \texttt{README.md}, \texttt{main.py},
\texttt{evidence.json}, and tests. The implementation-specific paths
are:

\begin{Shaded}
\begin{Highlighting}[]
\NormalTok{src/agent\_course/contracts.py}
\NormalTok{tests/unit/test\_contracts.py}
\end{Highlighting}
\end{Shaded}

These are \textbf{proposed student deliverables}, not claims that an
implemented repository was supplied.

\textbf{Implementation sequence:} (1) Define typed request/result/error
contracts; (2) implement Define ModelPort, ToolSpec, AgentRequest and
AgentResult; create a deterministic FakeModel that returns predefined
actions; never execute model-produced source code; (3) add fixtures and
logging with immutable run IDs; (4) write security-first negative tests
before any live integration; (5) commit a design/decision diagram; (6)
capture verification evidence and compare results with the week's stated
rubric.

\textbf{Positive test:} valid JSON plan with exactly one read-only tool
proposal.\\
\textbf{Negative/adversarial test:} unknown tool and invalid schema
rejected.\\
\textbf{Security requirement:} Model output is untrusted, and model
reasoning never authorizes an API call.\\
\textbf{Verification command} (target command after implementing the
listed test files):
\texttt{python\ -m\ pytest\ -q\ tests/unit/test\_contracts.py}.\\
\textbf{Troubleshooting:} If the model stub emits unexpected fields,
dump sanitized parsed keys and confirm schema version.

\textbf{Acceptance criteria:} All five layers correctly assigned; three
invalid inputs rejected; no tool invoked without validation; a clean
offline checkout must reproduce the result with a fixed test seed, and
\texttt{evidence.json} must record commit SHA, run ID, versions, check
status, and reviewer/tester. Reject a lab if any required
authorization-negative case fails.

\textbf{Assessment rubric (100 points):} correct architecture and
contracts 20; implementation and type correctness 25; positive/negative
tests 30; security and error handling 15; reproducibility and
documentation 10. \textbf{Submission:} source, fixture data, diagram,
tests, evidence log and a one-page engineering summary. \textbf{Advanced
extension:} port the contract to TypeScript/Java or run an independently
authorized cloud test; no advanced work can override failed safety
gates.

\subsubsection{L02 --- Minimal framework-free agent
execution}\label{l02-minimal-framework-free-agent-execution}

\textbf{Purpose and prerequisites:} Execute the Week 02 engineering
outcome with the prerequisites stated in Week 02. \textbf{Target time:}
4 h implementation plus 1.5 h verification (other course work handled in
the weekly allocation).

\textbf{Learning objective:} Develop a bounded reason--act--observe
state machine without an agent SDK. \textbf{Architecture diagram
(logical):}

\begin{Shaded}
\begin{Highlighting}[]
\NormalTok{Task → bounded controller → allowlisted dispatcher →}
\NormalTok{fake tool → event store}
\end{Highlighting}
\end{Shaded}

\textbf{Required tools:} Python 3, pytest, typed model stub, fixtures
and Git; add only the week-specific dependencies (MCP SDK in L06--07,
OPA in L10, OpenTelemetry in L13, Docker in L14); pin each one in the
lockfile.

\textbf{Starter files / repository structure:} create a directory
\texttt{labs/L02/} containing \texttt{README.md}, \texttt{main.py},
\texttt{evidence.json}, and tests. The implementation-specific paths
are:

\begin{Shaded}
\begin{Highlighting}[]
\NormalTok{harness/runner.py}
\NormalTok{harness/dispatcher.py}
\NormalTok{tests/unit/test\_runtime.py}
\end{Highlighting}
\end{Shaded}

These are \textbf{proposed student deliverables}, not claims that an
implemented repository was supplied.

\textbf{Implementation sequence:} (1) Define typed request/result/error
contracts; (2) implement Implement run\_id, max\_steps, deadline,
allowlisted tool catalog, approved input schema and append-only events;
(3) add fixtures and logging with immutable run IDs; (4) write
security-first negative tests before any live integration; (5) commit a
design/decision diagram; (6) capture verification evidence and compare
results with the week's stated rubric.

\textbf{Positive test:} read-only sequence terminates within the
configured step budget.\\
\textbf{Negative/adversarial test:} timeout, duplicate action, disabled
tool and exhausted steps blocked.\\
\textbf{Security requirement:} Only pure read-only tools enabled; all
write attempts denied.\\
\textbf{Verification command} (target command after implementing the
listed test files):
\texttt{python\ -m\ pytest\ -q\ tests/unit/test\_runtime.py}.\\
\textbf{Troubleshooting:} If tests hang, ensure deadline and
cancellation apply to both model and tool await points.

\textbf{Acceptance criteria:} No infinite loop, explicit reason for
every terminal state, duplicate invocation prevented; a clean offline
checkout must reproduce the result with a fixed test seed, and
\texttt{evidence.json} must record commit SHA, run ID, versions, check
status, and reviewer/tester. Reject a lab if any required
authorization-negative case fails.

\textbf{Assessment rubric (100 points):} correct architecture and
contracts 20; implementation and type correctness 25; positive/negative
tests 30; security and error handling 15; reproducibility and
documentation 10. \textbf{Submission:} source, fixture data, diagram,
tests, evidence log and a one-page engineering summary. \textbf{Advanced
extension:} port the contract to TypeScript/Java or run an independently
authorized cloud test; no advanced work can override failed safety
gates.

\subsubsection{L03 --- Framework orchestration: LangGraph and SDK
alternatives}\label{l03-framework-orchestration-langgraph-and-sdk-alternatives}

\textbf{Purpose and prerequisites:} Execute the Week 03 engineering
outcome with the prerequisites stated in Week 03. \textbf{Target time:}
4 h implementation plus 1.5 h verification (other course work handled in
the weekly allocation).

\textbf{Learning objective:} Implement equivalent routing using a graph
and the handcrafted runner. \textbf{Architecture diagram (logical):}

\begin{Shaded}
\begin{Highlighting}[]
\NormalTok{Typed state → LangGraph nodes → policy{-}backed tool}
\NormalTok{adapter → result}
\end{Highlighting}
\end{Shaded}

\textbf{Required tools:} Python 3, pytest, typed model stub, fixtures
and Git; add only the week-specific dependencies (MCP SDK in L06--07,
OPA in L10, OpenTelemetry in L13, Docker in L14); pin each one in the
lockfile.

\textbf{Starter files / repository structure:} create a directory
\texttt{labs/L03/} containing \texttt{README.md}, \texttt{main.py},
\texttt{evidence.json}, and tests. The implementation-specific paths
are:

\begin{Shaded}
\begin{Highlighting}[]
\NormalTok{framework/graph.py}
\NormalTok{docs/adrs/ADR{-}001.md}
\end{Highlighting}
\end{Shaded}

These are \textbf{proposed student deliverables}, not claims that an
implemented repository was supplied.

\textbf{Implementation sequence:} (1) Define typed request/result/error
contracts; (2) implement Build analysis→retrieve→verify→report nodes
with a typed state; maintain same external contracts as W2; (3) add
fixtures and logging with immutable run IDs; (4) write security-first
negative tests before any live integration; (5) commit a design/decision
diagram; (6) capture verification evidence and compare results with the
week's stated rubric.

\textbf{Positive test:} correct branch and report match the reference
implementation.\\
\textbf{Negative/adversarial test:} unauthorized graph branch and
malformed state rejected.\\
\textbf{Security requirement:} Enforce authorization in
application-owned adapters not only model instructions.\\
\textbf{Verification command} (target command after implementing the
listed test files):
\texttt{python\ -m\ pytest\ -q\ tests/contract/test\_graph\_parity.py}.\\
\textbf{Troubleshooting:} If state mutates unexpectedly, inspect reducer
merges and immutable fixture snapshots.

\textbf{Acceptance criteria:} Graph reproducible; trade-off evaluation
has ≥6 criteria and ≥3 technologies; a clean offline checkout must
reproduce the result with a fixed test seed, and \texttt{evidence.json}
must record commit SHA, run ID, versions, check status, and
reviewer/tester. Reject a lab if any required authorization-negative
case fails.

\textbf{Assessment rubric (100 points):} correct architecture and
contracts 20; implementation and type correctness 25; positive/negative
tests 30; security and error handling 15; reproducibility and
documentation 10. \textbf{Submission:} source, fixture data, diagram,
tests, evidence log and a one-page engineering summary. \textbf{Advanced
extension:} port the contract to TypeScript/Java or run an independently
authorized cloud test; no advanced work can override failed safety
gates.

\subsubsection{L04 --- Durable state, checkpoint recovery and human
interrupts}\label{l04-durable-state-checkpoint-recovery-and-human-interrupts}

\textbf{Purpose and prerequisites:} Execute the Week 04 engineering
outcome with the prerequisites stated in Week 04. \textbf{Target time:}
4 h implementation plus 1.5 h verification (other course work handled in
the weekly allocation).

\textbf{Learning objective:} Persist execution checkpoints with
monotonic state transitions. \textbf{Architecture diagram (logical):}

\begin{Shaded}
\begin{Highlighting}[]
\NormalTok{Workflow → checkpoint store ↔ approval gate → resume →}
\NormalTok{verifier}
\end{Highlighting}
\end{Shaded}

\textbf{Required tools:} Python 3, pytest, typed model stub, fixtures
and Git; add only the week-specific dependencies (MCP SDK in L06--07,
OPA in L10, OpenTelemetry in L13, Docker in L14); pin each one in the
lockfile.

\textbf{Starter files / repository structure:} create a directory
\texttt{labs/L04/} containing \texttt{README.md}, \texttt{main.py},
\texttt{evidence.json}, and tests. The implementation-specific paths
are:

\begin{Shaded}
\begin{Highlighting}[]
\NormalTok{harness/checkpoints.py}
\NormalTok{tests/integration/test\_resume.py}
\end{Highlighting}
\end{Shaded}

These are \textbf{proposed student deliverables}, not claims that an
implemented repository was supplied.

\textbf{Implementation sequence:} (1) Define typed request/result/error
contracts; (2) implement Persist state before/after each tool call and
link run\_id to trace\_id; implement pause for approval; (3) add
fixtures and logging with immutable run IDs; (4) write security-first
negative tests before any live integration; (5) commit a design/decision
diagram; (6) capture verification evidence and compare results with the
week's stated rubric.

\textbf{Positive test:} interrupted read resumes with exactly one final
report.\\
\textbf{Negative/adversarial test:} stale approval and restart after
partial operation prevented.\\
\textbf{Security requirement:} Approver identity and policy version
recorded; unapproved resume impossible.\\
\textbf{Verification command} (target command after implementing the
listed test files):
\texttt{python\ -m\ pytest\ -q\ tests/integration/test\_resume.py}.\\
\textbf{Troubleshooting:} If a replay repeats work, inspect idempotency
key, commit boundary and checkpoint ordering.

\textbf{Acceptance criteria:} Resumed run produces exactly one durable
outcome and preserves prior evidence; a clean offline checkout must
reproduce the result with a fixed test seed, and \texttt{evidence.json}
must record commit SHA, run ID, versions, check status, and
reviewer/tester. Reject a lab if any required authorization-negative
case fails.

\textbf{Assessment rubric (100 points):} correct architecture and
contracts 20; implementation and type correctness 25; positive/negative
tests 30; security and error handling 15; reproducibility and
documentation 10. \textbf{Submission:} source, fixture data, diagram,
tests, evidence log and a one-page engineering summary. \textbf{Advanced
extension:} port the contract to TypeScript/Java or run an independently
authorized cloud test; no advanced work can override failed safety
gates.

\subsubsection{L05 --- Agent harness control plane, context and
environment}\label{l05-agent-harness-control-plane-context-and-environment}

\textbf{Purpose and prerequisites:} Execute the Week 05 engineering
outcome with the prerequisites stated in Week 05. \textbf{Target time:}
4 h implementation plus 1.5 h verification (other course work handled in
the weekly allocation).

\textbf{Learning objective:} Separate tool exposure, invocation, policy
checks, and environment isolation. \textbf{Architecture diagram
(logical):}

\begin{Shaded}
\begin{Highlighting}[]
\NormalTok{Model ↔ harness control plane → policy gate → sandbox}
\NormalTok{adapter → event log}
\end{Highlighting}
\end{Shaded}

\textbf{Required tools:} Python 3, pytest, typed model stub, fixtures
and Git; add only the week-specific dependencies (MCP SDK in L06--07,
OPA in L10, OpenTelemetry in L13, Docker in L14); pin each one in the
lockfile.

\textbf{Starter files / repository structure:} create a directory
\texttt{labs/L05/} containing \texttt{README.md}, \texttt{main.py},
\texttt{evidence.json}, and tests. The implementation-specific paths
are:

\begin{Shaded}
\begin{Highlighting}[]
\NormalTok{harness/permissions.py}
\NormalTok{harness/budgets.py}
\NormalTok{sandbox config}
\end{Highlighting}
\end{Shaded}

These are \textbf{proposed student deliverables}, not claims that an
implemented repository was supplied.

\textbf{Implementation sequence:} (1) Define typed request/result/error
contracts; (2) implement Introduce policy gate, context store,
restricted execution adapter and explicit budget accounting to W2/W4
harness; (3) add fixtures and logging with immutable run IDs; (4) write
security-first negative tests before any live integration; (5) commit a
design/decision diagram; (6) capture verification evidence and compare
results with the week's stated rubric.

\textbf{Positive test:} authorized read succeeds under configured
execution budgets.\\
\textbf{Negative/adversarial test:} injected tool output and forbidden
filesystem/network action denied.\\
\textbf{Security requirement:} No shell/network by default; tool outputs
treated as untrusted.\\
\textbf{Verification command} (target command after implementing the
listed test files):
\texttt{python\ -m\ pytest\ -q\ tests/adversarial/test\_harness.py}.\\
\textbf{Troubleshooting:} If the sandbox leaks access, inspect adapter
allowlist and host-level container mounts.

\textbf{Acceptance criteria:} Enforcement remains effective if model
prompt is malicious or contradictory; a clean offline checkout must
reproduce the result with a fixed test seed, and \texttt{evidence.json}
must record commit SHA, run ID, versions, check status, and
reviewer/tester. Reject a lab if any required authorization-negative
case fails.

\textbf{Assessment rubric (100 points):} correct architecture and
contracts 20; implementation and type correctness 25; positive/negative
tests 30; security and error handling 15; reproducibility and
documentation 10. \textbf{Submission:} source, fixture data, diagram,
tests, evidence log and a one-page engineering summary. \textbf{Advanced
extension:} port the contract to TypeScript/Java or run an independently
authorized cloud test; no advanced work can override failed safety
gates.

\subsubsection{L06 --- MCP host/client/server fundamentals and protocol
revisions}\label{l06-mcp-hostclientserver-fundamentals-and-protocol-revisions}

\textbf{Purpose and prerequisites:} Execute the Week 06 engineering
outcome with the prerequisites stated in Week 06. \textbf{Target time:}
4 h implementation plus 1.5 h verification (other course work handled in
the weekly allocation).

\textbf{Learning objective:} Diagram exact host/client/server and
tool/result trust boundaries. \textbf{Architecture diagram (logical):}

\begin{Shaded}
\begin{Highlighting}[]
\NormalTok{Host [client] ⇄ MCP server [tools/resources/prompts] →}
\NormalTok{fixture store}
\end{Highlighting}
\end{Shaded}

\textbf{Required tools:} Python 3, pytest, typed model stub, fixtures
and Git; add only the week-specific dependencies (MCP SDK in L06--07,
OPA in L10, OpenTelemetry in L13, Docker in L14); pin each one in the
lockfile.

\textbf{Starter files / repository structure:} create a directory
\texttt{labs/L06/} containing \texttt{README.md}, \texttt{main.py},
\texttt{evidence.json}, and tests. The implementation-specific paths
are:

\begin{Shaded}
\begin{Highlighting}[]
\NormalTok{mcp/server.py}
\NormalTok{mcp/client.py}
\NormalTok{fixtures/mcp/*}
\end{Highlighting}
\end{Shaded}

These are \textbf{proposed student deliverables}, not claims that an
implemented repository was supplied.

\textbf{Implementation sequence:} (1) Define typed request/result/error
contracts; (2) implement Build mock issue-metadata MCP server with
read-only tools; run dual compatibility experiments against pinned SDK
versions; (3) add fixtures and logging with immutable run IDs; (4) write
security-first negative tests before any live integration; (5) commit a
design/decision diagram; (6) capture verification evidence and compare
results with the week's stated rubric.

\textbf{Positive test:} tools/list and tools/call return valid
schemas/results.\\
\textbf{Negative/adversarial test:} unsupported protocol revision,
forged tool and invalid resource ID rejected.\\
\textbf{Security requirement:} Do not expose entire conversation; deny
unexpected resource identifiers.\\
\textbf{Verification command} (target command after implementing the
listed test files):
\texttt{python\ -m\ pytest\ -q\ tests/contract/test\_mcp.py}.\\
\textbf{Troubleshooting:} If requests fail, compare protocol revision
and client/server transport rather than guessing API method names.

\textbf{Acceptance criteria:} Five MCP protocol cases documented;
correct capability behavior for chosen SDK version; a clean offline
checkout must reproduce the result with a fixed test seed, and
\texttt{evidence.json} must record commit SHA, run ID, versions, check
status, and reviewer/tester. Reject a lab if any required
authorization-negative case fails.

\textbf{Assessment rubric (100 points):} correct architecture and
contracts 20; implementation and type correctness 25; positive/negative
tests 30; security and error handling 15; reproducibility and
documentation 10. \textbf{Submission:} source, fixture data, diagram,
tests, evidence log and a one-page engineering summary. \textbf{Advanced
extension:} port the contract to TypeScript/Java or run an independently
authorized cloud test; no advanced work can override failed safety
gates.

\subsubsection{L07 --- Authenticated GitHub MCP and REST
integration}\label{l07-authenticated-github-mcp-and-rest-integration}

\textbf{Purpose and prerequisites:} Execute the Week 07 engineering
outcome with the prerequisites stated in Week 07. \textbf{Target time:}
4 h implementation plus 1.5 h verification (other course work handled in
the weekly allocation).

\textbf{Learning objective:} Design authorization boundaries for GitHub
App installation tokens. \textbf{Architecture diagram (logical):}

\begin{Shaded}
\begin{Highlighting}[]
\NormalTok{Harness → MCP client → scoped GitHub façade → fake REST}
\NormalTok{API → evidence}
\end{Highlighting}
\end{Shaded}

\textbf{Required tools:} Python 3, pytest, typed model stub, fixtures
and Git; add only the week-specific dependencies (MCP SDK in L06--07,
OPA in L10, OpenTelemetry in L13, Docker in L14); pin each one in the
lockfile.

\textbf{Starter files / repository structure:} create a directory
\texttt{labs/L07/} containing \texttt{README.md}, \texttt{main.py},
\texttt{evidence.json}, and tests. The implementation-specific paths
are:

\begin{Shaded}
\begin{Highlighting}[]
\NormalTok{mcp/github\_gateway.py}
\NormalTok{fixtures/github/*}
\end{Highlighting}
\end{Shaded}

These are \textbf{proposed student deliverables}, not claims that an
implemented repository was supplied.

\textbf{Implementation sequence:} (1) Define typed request/result/error
contracts; (2) implement Create get\_repo\_metadata, list\_alerts,
get\_alert\_evidence contracts; mock GitHub API first; optional verified
test installation; (3) add fixtures and logging with immutable run IDs;
(4) write security-first negative tests before any live integration; (5)
commit a design/decision diagram; (6) capture verification evidence and
compare results with the week's stated rubric.

\textbf{Positive test:} authorized repository alerts paginated and
normalized.\\
\textbf{Negative/adversarial test:} cross-repository access, expired
token, rate-limit failure handled safely.\\
\textbf{Security requirement:} Separate read and write identities; no
broad personal access tokens in committed configs.\\
\textbf{Verification command} (target command after implementing the
listed test files):
\texttt{python\ -m\ pytest\ -q\ tests/contract/test\_github.py}.\\
\textbf{Troubleshooting:} If a request returns 403, inspect installation
scope and repository permissions, not prompt wording.

\textbf{Acceptance criteria:} 0 successful forbidden writes; paginated
alerts and 429 handled without data loss; a clean offline checkout must
reproduce the result with a fixed test seed, and \texttt{evidence.json}
must record commit SHA, run ID, versions, check status, and
reviewer/tester. Reject a lab if any required authorization-negative
case fails.

\textbf{Assessment rubric (100 points):} correct architecture and
contracts 20; implementation and type correctness 25; positive/negative
tests 30; security and error handling 15; reproducibility and
documentation 10. \textbf{Submission:} source, fixture data, diagram,
tests, evidence log and a one-page engineering summary. \textbf{Advanced
extension:} port the contract to TypeScript/Java or run an independently
authorized cloud test; no advanced work can override failed safety
gates.

\subsubsection{L08 --- Context engineering, retrieval and midterm
demonstration}\label{l08-context-engineering-retrieval-and-midterm-demonstration}

\textbf{Purpose and prerequisites:} Execute the Week 08 engineering
outcome with the prerequisites stated in Week 08. \textbf{Target time:}
4 h implementation plus 1.5 h verification (other course work handled in
the weekly allocation).

\textbf{Learning objective:} Compare lexical and semantic retrieval
quality with a labeled corpus. \textbf{Architecture diagram (logical):}

\begin{Shaded}
\begin{Highlighting}[]
\NormalTok{Query → retrieval index → context filter → harness →}
\NormalTok{cited report}
\end{Highlighting}
\end{Shaded}

\textbf{Required tools:} Python 3, pytest, typed model stub, fixtures
and Git; add only the week-specific dependencies (MCP SDK in L06--07,
OPA in L10, OpenTelemetry in L13, Docker in L14); pin each one in the
lockfile.

\textbf{Starter files / repository structure:} create a directory
\texttt{labs/L08/} containing \texttt{README.md}, \texttt{main.py},
\texttt{evidence.json}, and tests. The implementation-specific paths
are:

\begin{Shaded}
\begin{Highlighting}[]
\NormalTok{retrieval/index.py}
\NormalTok{retrieval/evidence.py}
\NormalTok{fixtures/eval\_queries.jsonl}
\end{Highlighting}
\end{Shaded}

These are \textbf{proposed student deliverables}, not claims that an
implemented repository was supplied.

\textbf{Implementation sequence:} (1) Define typed request/result/error
contracts; (2) implement Index 50 synthetic alerts and code excerpts;
retrieve exact commit+rule evidence and produce citation links; (3) add
fixtures and logging with immutable run IDs; (4) write security-first
negative tests before any live integration; (5) commit a design/decision
diagram; (6) capture verification evidence and compare results with the
week's stated rubric.

\textbf{Positive test:} relevant evidence returned with resolvable
citations.\\
\textbf{Negative/adversarial test:} poisoned retrieval chunk and missing
commit provenance escalated.\\
\textbf{Security requirement:} Untrusted retrieved text never changes
policy or tool catalog.\\
\textbf{Verification command} (target command after implementing the
listed test files):
\texttt{python\ -m\ pytest\ -q\ tests/evaluation/test\_retrieval.py}.\\
\textbf{Troubleshooting:} If retrieval degrades, inspect chunk overlap,
index freshness and golden labels.

\textbf{Acceptance criteria:} All cited artifacts resolve to fixture
IDs; precision@3 measured; no unsupported permission change; a clean
offline checkout must reproduce the result with a fixed test seed, and
\texttt{evidence.json} must record commit SHA, run ID, versions, check
status, and reviewer/tester. Reject a lab if any required
authorization-negative case fails.

\textbf{Assessment rubric (100 points):} correct architecture and
contracts 20; implementation and type correctness 25; positive/negative
tests 30; security and error handling 15; reproducibility and
documentation 10. \textbf{Submission:} source, fixture data, diagram,
tests, evidence log and a one-page engineering summary. \textbf{Advanced
extension:} port the contract to TypeScript/Java or run an independently
authorized cloud test; no advanced work can override failed safety
gates.

\subsubsection{L09 --- Agent threat modeling and adversarial
boundaries}\label{l09-agent-threat-modeling-and-adversarial-boundaries}

\textbf{Purpose and prerequisites:} Execute the Week 09 engineering
outcome with the prerequisites stated in Week 09. \textbf{Target time:}
4 h implementation plus 1.5 h verification (other course work handled in
the weekly allocation).

\textbf{Learning objective:} Construct STRIDE-informed abuse cases
against model/tool/MCP boundaries. \textbf{Architecture diagram
(logical):}

\begin{Shaded}
\begin{Highlighting}[]
\NormalTok{Untrusted issue/tool output → harness PEP → policy}
\NormalTok{decision → deny/report}
\end{Highlighting}
\end{Shaded}

\textbf{Required tools:} Python 3, pytest, typed model stub, fixtures
and Git; add only the week-specific dependencies (MCP SDK in L06--07,
OPA in L10, OpenTelemetry in L13, Docker in L14); pin each one in the
lockfile.

\textbf{Starter files / repository structure:} create a directory
\texttt{labs/L09/} containing \texttt{README.md}, \texttt{main.py},
\texttt{evidence.json}, and tests. The implementation-specific paths
are:

\begin{Shaded}
\begin{Highlighting}[]
\NormalTok{docs/threat{-}models/agent{-}threats.md}
\NormalTok{fixtures/attack\_cases/*}
\end{Highlighting}
\end{Shaded}

These are \textbf{proposed student deliverables}, not claims that an
implemented repository was supplied.

\textbf{Implementation sequence:} (1) Define typed request/result/error
contracts; (2) implement Create 12 attacks: forged system directions in
issue body, malicious MCP tool description, fake approval, tainted
retrieved snippet; (3) add fixtures and logging with immutable run IDs;
(4) write security-first negative tests before any live integration; (5)
commit a design/decision diagram; (6) capture verification evidence and
compare results with the week's stated rubric.

\textbf{Positive test:} ordinary issue text is summarized without
affecting policy.\\
\textbf{Negative/adversarial test:} 12 malicious prompts/tool
descriptors cannot trigger write.\\
\textbf{Security requirement:} Assume repo, issues and tool outputs are
attacker-controlled data.\\
\textbf{Verification command} (target command after implementing the
listed test files):
\texttt{python\ -m\ pytest\ -q\ tests/adversarial/test\_injection.py}.\\
\textbf{Troubleshooting:} If injection appears effective, isolate
trusted instruction channel from retrieved content.

\textbf{Acceptance criteria:} ≥12 attack cases; 100\% of
authorization-bypass fixtures blocked; a clean offline checkout must
reproduce the result with a fixed test seed, and \texttt{evidence.json}
must record commit SHA, run ID, versions, check status, and
reviewer/tester. Reject a lab if any required authorization-negative
case fails.

\textbf{Assessment rubric (100 points):} correct architecture and
contracts 20; implementation and type correctness 25; positive/negative
tests 30; security and error handling 15; reproducibility and
documentation 10. \textbf{Submission:} source, fixture data, diagram,
tests, evidence log and a one-page engineering summary. \textbf{Advanced
extension:} port the contract to TypeScript/Java or run an independently
authorized cloud test; no advanced work can override failed safety
gates.

\subsubsection{L10 --- Policy-as-code, human approval and safe
actions}\label{l10-policy-as-code-human-approval-and-safe-actions}

\textbf{Purpose and prerequisites:} Execute the Week 10 engineering
outcome with the prerequisites stated in Week 10. \textbf{Target time:}
4 h implementation plus 1.5 h verification (other course work handled in
the weekly allocation).

\textbf{Learning objective:} Separate AI recommendation from
deterministic policy adjudication. \textbf{Architecture diagram
(logical):}

\begin{Shaded}
\begin{Highlighting}[]
\NormalTok{Model proposal → approval queue → OPA PDP → mock write}
\NormalTok{PEP → audit}
\end{Highlighting}
\end{Shaded}

\textbf{Required tools:} Python 3, pytest, typed model stub, fixtures
and Git; add only the week-specific dependencies (MCP SDK in L06--07,
OPA in L10, OpenTelemetry in L13, Docker in L14); pin each one in the
lockfile.

\textbf{Starter files / repository structure:} create a directory
\texttt{labs/L10/} containing \texttt{README.md}, \texttt{main.py},
\texttt{evidence.json}, and tests. The implementation-specific paths
are:

\begin{Shaded}
\begin{Highlighting}[]
\NormalTok{policy/agent.rego}
\NormalTok{policy/agent\_test.rego}
\NormalTok{mock action adapter}
\end{Highlighting}
\end{Shaded}

These are \textbf{proposed student deliverables}, not claims that an
implemented repository was supplied.

\textbf{Implementation sequence:} (1) Define typed request/result/error
contracts; (2) implement Create proposed-disposition schema and OPA
policy; stage a mock dismiss-alert request but never alter a live alert;
(3) add fixtures and logging with immutable run IDs; (4) write
security-first negative tests before any live integration; (5) commit a
design/decision diagram; (6) capture verification evidence and compare
results with the week's stated rubric.

\textbf{Positive test:} mock approval accepted with matching identity,
scope and policy version.\\
\textbf{Negative/adversarial test:} expired, revoked and replayed
approvals denied without side effect.\\
\textbf{Security requirement:} Sensitive write operations restricted to
a separate mocked service behind authenticated approval.\\
\textbf{Verification command} (target command after implementing the
listed test files):
\texttt{opa\ test\ policy/\ -v\ \&\&\ python\ -m\ pytest\ -q\ tests/adversarial/test\_approvals.py}.\\
\textbf{Troubleshooting:} If policy unexpectedly allows, test undefined
input and explicit default deny.

\textbf{Acceptance criteria:} All negative gate tests pass; no
side-effect if approval absent; a clean offline checkout must reproduce
the result with a fixed test seed, and \texttt{evidence.json} must
record commit SHA, run ID, versions, check status, and reviewer/tester.
Reject a lab if any required authorization-negative case fails.

\textbf{Assessment rubric (100 points):} correct architecture and
contracts 20; implementation and type correctness 25; positive/negative
tests 30; security and error handling 15; reproducibility and
documentation 10. \textbf{Submission:} source, fixture data, diagram,
tests, evidence log and a one-page engineering summary. \textbf{Advanced
extension:} port the contract to TypeScript/Java or run an independently
authorized cloud test; no advanced work can override failed safety
gates.

\subsubsection{L11 --- Multi-agent delegation and distributed
execution}\label{l11-multi-agent-delegation-and-distributed-execution}

\textbf{Purpose and prerequisites:} Execute the Week 11 engineering
outcome with the prerequisites stated in Week 11. \textbf{Target time:}
4 h implementation plus 1.5 h verification (other course work handled in
the weekly allocation).

\textbf{Learning objective:} Implement supervisor, investigator, critic
and verifier roles. \textbf{Architecture diagram (logical):}

\begin{Shaded}
\begin{Highlighting}[]
\NormalTok{Supervisor → isolated investigator/critic → verifier →}
\NormalTok{evidence ledger}
\end{Highlighting}
\end{Shaded}

\textbf{Required tools:} Python 3, pytest, typed model stub, fixtures
and Git; add only the week-specific dependencies (MCP SDK in L06--07,
OPA in L10, OpenTelemetry in L13, Docker in L14); pin each one in the
lockfile.

\textbf{Starter files / repository structure:} create a directory
\texttt{labs/L11/} containing \texttt{README.md}, \texttt{main.py},
\texttt{evidence.json}, and tests. The implementation-specific paths
are:

\begin{Shaded}
\begin{Highlighting}[]
\NormalTok{framework/delegation.py}
\NormalTok{tests/integration/test\_delegation.py}
\end{Highlighting}
\end{Shaded}

These are \textbf{proposed student deliverables}, not claims that an
implemented repository was supplied.

\textbf{Implementation sequence:} (1) Define typed request/result/error
contracts; (2) implement Run planner→investigator→reviewer with
read-only tools and immutable evidence snapshots; verifier owns
completion criteria; (3) add fixtures and logging with immutable run
IDs; (4) write security-first negative tests before any live
integration; (5) commit a design/decision diagram; (6) capture
verification evidence and compare results with the week's stated rubric.

\textbf{Positive test:} reviewer verifies a delegated task with bounded
concurrency.\\
\textbf{Negative/adversarial test:} worker failure, task cancellation
and delegation-depth overflow controlled.\\
\textbf{Security requirement:} Each worker receives narrowed credentials
and tool set; supervisor cannot create grants.\\
\textbf{Verification command} (target command after implementing the
listed test files):
\texttt{python\ -m\ pytest\ -q\ tests/integration/test\_delegation.py}.\\
\textbf{Troubleshooting:} If tasks deadlock, inspect cancellation
propagation and bounded worker queue.

\textbf{Acceptance criteria:} No worker privilege escalation; replay
demonstrates deterministic verification; a clean offline checkout must
reproduce the result with a fixed test seed, and \texttt{evidence.json}
must record commit SHA, run ID, versions, check status, and
reviewer/tester. Reject a lab if any required authorization-negative
case fails.

\textbf{Assessment rubric (100 points):} correct architecture and
contracts 20; implementation and type correctness 25; positive/negative
tests 30; security and error handling 15; reproducibility and
documentation 10. \textbf{Submission:} source, fixture data, diagram,
tests, evidence log and a one-page engineering summary. \textbf{Advanced
extension:} port the contract to TypeScript/Java or run an independently
authorized cloud test; no advanced work can override failed safety
gates.

\subsubsection{L12 --- Evaluation-driven agent quality
engineering}\label{l12-evaluation-driven-agent-quality-engineering}

\textbf{Purpose and prerequisites:} Execute the Week 12 engineering
outcome with the prerequisites stated in Week 12. \textbf{Target time:}
4 h implementation plus 1.5 h verification (other course work handled in
the weekly allocation).

\textbf{Learning objective:} Build labeled datasets and distinguish
deterministic checks from judge-model ratings. \textbf{Architecture
diagram (logical):}

\begin{Shaded}
\begin{Highlighting}[]
\NormalTok{Golden fixtures → agent pipeline → evaluator → metrics}
\NormalTok{→ CI gate}
\end{Highlighting}
\end{Shaded}

\textbf{Required tools:} Python 3, pytest, typed model stub, fixtures
and Git; add only the week-specific dependencies (MCP SDK in L06--07,
OPA in L10, OpenTelemetry in L13, Docker in L14); pin each one in the
lockfile.

\textbf{Starter files / repository structure:} create a directory
\texttt{labs/L12/} containing \texttt{README.md}, \texttt{main.py},
\texttt{evidence.json}, and tests. The implementation-specific paths
are:

\begin{Shaded}
\begin{Highlighting}[]
\NormalTok{tests/evaluation/gold.jsonl}
\NormalTok{scripts/eval\_gate.py}
\end{Highlighting}
\end{Shaded}

These are \textbf{proposed student deliverables}, not claims that an
implemented repository was supplied.

\textbf{Implementation sequence:} (1) Define typed request/result/error
contracts; (2) implement Create 100 synthetic alerts: 40 legitimate, 40
benign candidates, 20 indeterminate; evaluate only recommendations and
evidence; (3) add fixtures and logging with immutable run IDs; (4) write
security-first negative tests before any live integration; (5) commit a
design/decision diagram; (6) capture verification evidence and compare
results with the week's stated rubric.

\textbf{Positive test:} seeded evaluation suite emits reproducible
metrics.\\
\textbf{Negative/adversarial test:} unsupported dismissal recommendation
fails evaluator.\\
\textbf{Security requirement:} Any unsupported dismissal suggestion or
denied operation is a zero-tolerance event in release gate.\\
\textbf{Verification command} (target command after implementing the
listed test files):
\texttt{python\ -m\ pytest\ -q\ tests/evaluation/\ \&\&\ python\ scripts/eval\_gate.py}.\\
\textbf{Troubleshooting:} If metrics drift, verify fixture immutability,
random seed and evaluator revision.

\textbf{Acceptance criteria:} Deterministic policy tests green; metrics
reproducible on fixed fixtures; a clean offline checkout must reproduce
the result with a fixed test seed, and \texttt{evidence.json} must
record commit SHA, run ID, versions, check status, and reviewer/tester.
Reject a lab if any required authorization-negative case fails.

\textbf{Assessment rubric (100 points):} correct architecture and
contracts 20; implementation and type correctness 25; positive/negative
tests 30; security and error handling 15; reproducibility and
documentation 10. \textbf{Submission:} source, fixture data, diagram,
tests, evidence log and a one-page engineering summary. \textbf{Advanced
extension:} port the contract to TypeScript/Java or run an independently
authorized cloud test; no advanced work can override failed safety
gates.

\subsubsection{L13 --- Observability, SLOs and operational fault
recovery}\label{l13-observability-slos-and-operational-fault-recovery}

\textbf{Purpose and prerequisites:} Execute the Week 13 engineering
outcome with the prerequisites stated in Week 13. \textbf{Target time:}
4 h implementation plus 1.5 h verification (other course work handled in
the weekly allocation).

\textbf{Learning objective:} Correlate model, harness, MCP, policy and
GitHub operation spans. \textbf{Architecture diagram (logical):}

\begin{Shaded}
\begin{Highlighting}[]
\NormalTok{Agent spans → OTel collector → metrics/traces ↔}
\NormalTok{incident tests}
\end{Highlighting}
\end{Shaded}

\textbf{Required tools:} Python 3, pytest, typed model stub, fixtures
and Git; add only the week-specific dependencies (MCP SDK in L06--07,
OPA in L10, OpenTelemetry in L13, Docker in L14); pin each one in the
lockfile.

\textbf{Starter files / repository structure:} create a directory
\texttt{labs/L13/} containing \texttt{README.md}, \texttt{main.py},
\texttt{evidence.json}, and tests. The implementation-specific paths
are:

\begin{Shaded}
\begin{Highlighting}[]
\NormalTok{telemetry/spans.py}
\NormalTok{docs/runbooks/incidents.md}
\end{Highlighting}
\end{Shaded}

These are \textbf{proposed student deliverables}, not claims that an
implemented repository was supplied.

\textbf{Implementation sequence:} (1) Define typed request/result/error
contracts; (2) implement Emit trace\_id/run\_id/tool\_call\_id; expose
metrics for latency, cost, denial and replay outcomes; define SLOs; (3)
add fixtures and logging with immutable run IDs; (4) write
security-first negative tests before any live integration; (5) commit a
design/decision diagram; (6) capture verification evidence and compare
results with the week's stated rubric.

\textbf{Positive test:} each accepted run emits correlated spans and
metrics.\\
\textbf{Negative/adversarial test:} secret in span and double-delivery
fail safely.\\
\textbf{Security requirement:} Redact source and tokens; keep immutable
audit references separate from raw prompts.\\
\textbf{Verification command} (target command after implementing the
listed test files):
\texttt{python\ -m\ pytest\ -q\ tests/integration/test\_tracing.py}.\\
\textbf{Troubleshooting:} If correlation breaks, propagate
run\_id/traceparent across MCP boundary.

\textbf{Acceptance criteria:} Every sampled run traces to policy
decision; redaction tests pass; a clean offline checkout must reproduce
the result with a fixed test seed, and \texttt{evidence.json} must
record commit SHA, run ID, versions, check status, and reviewer/tester.
Reject a lab if any required authorization-negative case fails.

\textbf{Assessment rubric (100 points):} correct architecture and
contracts 20; implementation and type correctness 25; positive/negative
tests 30; security and error handling 15; reproducibility and
documentation 10. \textbf{Submission:} source, fixture data, diagram,
tests, evidence log and a one-page engineering summary. \textbf{Advanced
extension:} port the contract to TypeScript/Java or run an independently
authorized cloud test; no advanced work can override failed safety
gates.

\subsubsection{L14 --- Containerized deployment, CI/CD and workload
isolation}\label{l14-containerized-deployment-cicd-and-workload-isolation}

\textbf{Purpose and prerequisites:} Execute the Week 14 engineering
outcome with the prerequisites stated in Week 14. \textbf{Target time:}
4 h implementation plus 1.5 h verification (other course work handled in
the weekly allocation).

\textbf{Learning objective:} Package and scan a containerized harness
service. \textbf{Architecture diagram (logical):}

\begin{Shaded}
\begin{Highlighting}[]
\NormalTok{Source → lint/tests → policy/scan gates → container →}
\NormalTok{sandbox runtime}
\end{Highlighting}
\end{Shaded}

\textbf{Required tools:} Python 3, pytest, typed model stub, fixtures
and Git; add only the week-specific dependencies (MCP SDK in L06--07,
OPA in L10, OpenTelemetry in L13, Docker in L14); pin each one in the
lockfile.

\textbf{Starter files / repository structure:} create a directory
\texttt{labs/L14/} containing \texttt{README.md}, \texttt{main.py},
\texttt{evidence.json}, and tests. The implementation-specific paths
are:

\begin{Shaded}
\begin{Highlighting}[]
\NormalTok{Dockerfile}
\NormalTok{.github/workflows/agent{-}course{-}ci.yml}
\NormalTok{release manifest}
\end{Highlighting}
\end{Shaded}

These are \textbf{proposed student deliverables}, not claims that an
implemented repository was supplied.

\textbf{Implementation sequence:} (1) Define typed request/result/error
contracts; (2) implement Build local Docker Compose demo with mock
providers; optional AWS deployment is designed but not provisioned; (3)
add fixtures and logging with immutable run IDs; (4) write
security-first negative tests before any live integration; (5) commit a
design/decision diagram; (6) capture verification evidence and compare
results with the week's stated rubric.

\textbf{Positive test:} scanned non-root container runs mock integration
smoke test.\\
\textbf{Negative/adversarial test:} missing lockfile, vulnerable
dependency and prohibited permission fail gate.\\
\textbf{Security requirement:} No secrets in logs/images; nonroot image;
no automatic production write scope.\\
\textbf{Verification command} (target command after implementing the
listed test files):
\texttt{docker\ build\ -t\ agent-course:test\ .\ \&\&\ python\ -m\ pytest\ -q\ tests/integration/test\_container.py}.\\
\textbf{Troubleshooting:} If a container cannot run, verify UID, volume
permissions, pinned dependencies and architecture.

\textbf{Acceptance criteria:} All gates pass in local/dry-run CI;
deployment design documents rollback; a clean offline checkout must
reproduce the result with a fixed test seed, and \texttt{evidence.json}
must record commit SHA, run ID, versions, check status, and
reviewer/tester. Reject a lab if any required authorization-negative
case fails.

\textbf{Assessment rubric (100 points):} correct architecture and
contracts 20; implementation and type correctness 25; positive/negative
tests 30; security and error handling 15; reproducibility and
documentation 10. \textbf{Submission:} source, fixture data, diagram,
tests, evidence log and a one-page engineering summary. \textbf{Advanced
extension:} port the contract to TypeScript/Java or run an independently
authorized cloud test; no advanced work can override failed safety
gates.

\subsubsection{L15 --- Capstone integration: GitHub AppSec agent
(read-only)}\label{l15-capstone-integration-github-appsec-agent-read-only}

\textbf{Purpose and prerequisites:} Execute the Week 15 engineering
outcome with the prerequisites stated in Week 15. \textbf{Target time:}
4 h implementation plus 1.5 h verification (other course work handled in
the weekly allocation).

\textbf{Learning objective:} Integrate tools, harness, MCP, retrieval,
policy and tracing into one system. \textbf{Architecture diagram
(logical):}

\begin{Shaded}
\begin{Highlighting}[]
\NormalTok{GitHub mock/API → MCP client → harness → evidence}
\NormalTok{verifier → report}
\end{Highlighting}
\end{Shaded}

\textbf{Required tools:} Python 3, pytest, typed model stub, fixtures
and Git; add only the week-specific dependencies (MCP SDK in L06--07,
OPA in L10, OpenTelemetry in L13, Docker in L14); pin each one in the
lockfile.

\textbf{Starter files / repository structure:} create a directory
\texttt{labs/L15/} containing \texttt{README.md}, \texttt{main.py},
\texttt{evidence.json}, and tests. The implementation-specific paths
are:

\begin{Shaded}
\begin{Highlighting}[]
\NormalTok{src/agent\_course/appsec/investigator.py}
\NormalTok{docs/evidence/*}
\end{Highlighting}
\end{Shaded}

These are \textbf{proposed student deliverables}, not claims that an
implemented repository was supplied.

\textbf{Implementation sequence:} (1) Define typed request/result/error
contracts; (2) implement Implement intake→fetch→analyze→verify→report;
separate optional mocked approval service; (3) add fixtures and logging
with immutable run IDs; (4) write security-first negative tests before
any live integration; (5) commit a design/decision diagram; (6) capture
verification evidence and compare results with the week's stated rubric.

\textbf{Positive test:} three distinct CodeQL fixtures produce sourced
outcomes.\\
\textbf{Negative/adversarial test:} missing evidence and attempted write
show blocked outcomes.\\
\textbf{Security requirement:} Forbidden write APIs not registered in
production read-only tool catalog.\\
\textbf{Verification command} (target command after implementing the
listed test files):
\texttt{python\ -m\ pytest\ -q\ tests/integration/test\_capstone.py}.\\
\textbf{Troubleshooting:} If evidence references break, use
content-addressed snapshots and revision-specific URLs.

\textbf{Acceptance criteria:} All recommendations reproducible and
source-backed; zero live writes; a clean offline checkout must reproduce
the result with a fixed test seed, and \texttt{evidence.json} must
record commit SHA, run ID, versions, check status, and reviewer/tester.
Reject a lab if any required authorization-negative case fails.

\textbf{Assessment rubric (100 points):} correct architecture and
contracts 20; implementation and type correctness 25; positive/negative
tests 30; security and error handling 15; reproducibility and
documentation 10. \textbf{Submission:} source, fixture data, diagram,
tests, evidence log and a one-page engineering summary. \textbf{Advanced
extension:} port the contract to TypeScript/Java or run an independently
authorized cloud test; no advanced work can override failed safety
gates.

\subsubsection{L16 --- Capstone hardening, independent review and
operational
handoff}\label{l16-capstone-hardening-independent-review-and-operational-handoff}

\textbf{Purpose and prerequisites:} Execute the Week 16 engineering
outcome with the prerequisites stated in Week 16. \textbf{Target time:}
4 h implementation plus 1.5 h verification (other course work handled in
the weekly allocation).

\textbf{Learning objective:} Conduct architecture, adversarial and
recovery reviews against acceptance criteria. \textbf{Architecture
diagram (logical):}

\begin{Shaded}
\begin{Highlighting}[]
\NormalTok{Frozen commit + fixtures → test/replay harness → review}
\NormalTok{committee → decision}
\end{Highlighting}
\end{Shaded}

\textbf{Required tools:} Python 3, pytest, typed model stub, fixtures
and Git; add only the week-specific dependencies (MCP SDK in L06--07,
OPA in L10, OpenTelemetry in L13, Docker in L14); pin each one in the
lockfile.

\textbf{Starter files / repository structure:} create a directory
\texttt{labs/L16/} containing \texttt{README.md}, \texttt{main.py},
\texttt{evidence.json}, and tests. The implementation-specific paths
are:

\begin{Shaded}
\begin{Highlighting}[]
\NormalTok{docs/reviews/readiness.md}
\NormalTok{docs/evidence/final/*}
\end{Highlighting}
\end{Shaded}

These are \textbf{proposed student deliverables}, not claims that an
implemented repository was supplied.

\textbf{Implementation sequence:} (1) Define typed request/result/error
contracts; (2) implement Freeze version manifest and evidence bundle;
run safety replay plus user-acceptance demonstrations; (3) add fixtures
and logging with immutable run IDs; (4) write security-first negative
tests before any live integration; (5) commit a design/decision diagram;
(6) capture verification evidence and compare results with the week's
stated rubric.

\textbf{Positive test:} frozen release can be reproduced and
independently reviewed.\\
\textbf{Negative/adversarial test:} 403, injection, crash, replay and
missing approval yield failed/blocked release.\\
\textbf{Security requirement:} No live dismissal; all state-changing
methods remain mocked until separately approved.\\
\textbf{Verification command} (target command after implementing the
listed test files):
\texttt{make\ lint\ test\ verify\ \&\&\ python\ scripts/review\_packet.py}.\\
\textbf{Troubleshooting:} If review fails to reproduce, compare
dependency lock, fixture hashes and CI artifact provenance.

\textbf{Acceptance criteria:} Safety gates all pass; reviewer can
reproduce build, evaluation and evidence from pinned manifest; a clean
offline checkout must reproduce the result with a fixed test seed, and
\texttt{evidence.json} must record commit SHA, run ID, versions, check
status, and reviewer/tester. Reject a lab if any required
authorization-negative case fails.

\textbf{Assessment rubric (100 points):} correct architecture and
contracts 20; implementation and type correctness 25; positive/negative
tests 30; security and error handling 15; reproducibility and
documentation 10. \textbf{Submission:} source, fixture data, diagram,
tests, evidence log and a one-page engineering summary. \textbf{Advanced
extension:} port the contract to TypeScript/Java or run an independently
authorized cloud test; no advanced work can override failed safety
gates.

\subsubsection{Example starter-code requirement: a bounded
runner}\label{example-starter-code-requirement-a-bounded-runner}

For L02, the following small \emph{illustrative} interface makes the
loop contract concrete. The student's implementation adds checkpoint
persistence, cancellation, outcome verification and structured event
logging; this sketch is not a production-ready harness and intentionally
exposes no unrestricted shell or network execution.

\begin{Shaded}
\begin{Highlighting}[]
\ImportTok{from}\NormalTok{ dataclasses }\ImportTok{import}\NormalTok{ dataclass}
\ImportTok{from}\NormalTok{ typing }\ImportTok{import}\NormalTok{ Any, Callable}

\AttributeTok{@dataclass}\NormalTok{(frozen}\OperatorTok{=}\VariableTok{True}\NormalTok{)}
\KeywordTok{class}\NormalTok{ Proposal:}
\NormalTok{    tool\_name: }\BuiltInTok{str}
\NormalTok{    arguments: }\BuiltInTok{dict}\NormalTok{[}\BuiltInTok{str}\NormalTok{, Any]}

\AttributeTok{@dataclass}\NormalTok{(frozen}\OperatorTok{=}\VariableTok{True}\NormalTok{)}
\KeywordTok{class}\NormalTok{ RunResult:}
\NormalTok{    status: }\BuiltInTok{str}
\NormalTok{    observations: }\BuiltInTok{tuple}\NormalTok{[Any, ...]}

\KeywordTok{class}\NormalTok{ BudgetExhausted(}\PreprocessorTok{Exception}\NormalTok{):}
    \ControlFlowTok{pass}

\KeywordTok{def}\NormalTok{ execute\_bounded(}
\NormalTok{    propose: Callable[[}\BuiltInTok{tuple}\NormalTok{[Any, ...]], Proposal }\OperatorTok{|} \VariableTok{None}\NormalTok{],}
\NormalTok{    dispatch: }\BuiltInTok{dict}\NormalTok{[}\BuiltInTok{str}\NormalTok{, Callable[..., Any]],}
    \OperatorTok{*}\NormalTok{, max\_steps: }\BuiltInTok{int} \OperatorTok{=} \DecValTok{5}\NormalTok{,}
\NormalTok{) }\OperatorTok{{-}\textgreater{}}\NormalTok{ RunResult:}
\NormalTok{    history: }\BuiltInTok{tuple}\NormalTok{[Any, ...] }\OperatorTok{=}\NormalTok{ ()}
    \ControlFlowTok{for}\NormalTok{ \_ }\KeywordTok{in} \BuiltInTok{range}\NormalTok{(max\_steps):}
\NormalTok{        proposal }\OperatorTok{=}\NormalTok{ propose(history)}
        \ControlFlowTok{if}\NormalTok{ proposal }\KeywordTok{is} \VariableTok{None}\NormalTok{:}
            \ControlFlowTok{return}\NormalTok{ RunResult(}\StringTok{"completed"}\NormalTok{, history)}
        \CommentTok{\# The per{-}tool argument schemas and authorization checks must run here.}
        \ControlFlowTok{if}\NormalTok{ proposal.tool\_name }\KeywordTok{not} \KeywordTok{in}\NormalTok{ dispatch:}
            \ControlFlowTok{return}\NormalTok{ RunResult(}\StringTok{"blocked"}\NormalTok{, history)}
\NormalTok{        history }\OperatorTok{+=}\NormalTok{ (dispatch[proposal.tool\_name](}\OperatorTok{**}\NormalTok{proposal.arguments),)}
    \ControlFlowTok{return}\NormalTok{ RunResult(}\StringTok{"budget\_exhausted"}\NormalTok{, history)}
\end{Highlighting}
\end{Shaded}

\textbf{Critical gap in the example:} a real harness must validate each
tool's arguments, enforce the caller's authorization and run budget,
control side effects, use timeouts/cancellation, redact events and
verify successful completion; \texttt{dispatch} membership alone is not
sufficient. Tests must prove all of these before treating the runner as
production-capable.

\subsection{14. Software Architecture and Documentation
Requirements}\label{software-architecture-and-documentation-requirements}

\begin{landscape}
\begin{longtable}[]{@{}
  >{\raggedright\arraybackslash}p{(\columnwidth - 6\tabcolsep) * \real{0.1900}}
  >{\raggedright\arraybackslash}p{(\columnwidth - 6\tabcolsep) * \real{0.4000}}
  >{\raggedright\arraybackslash}p{(\columnwidth - 6\tabcolsep) * \real{0.1400}}
  >{\raggedright\arraybackslash}p{(\columnwidth - 6\tabcolsep) * \real{0.2700}}@{}}
\toprule\noalign{}
\begin{minipage}[b]{\linewidth}\raggedright
Artifact
\end{minipage} & \begin{minipage}[b]{\linewidth}\raggedright
Minimum contents
\end{minipage} & \begin{minipage}[b]{\linewidth}\raggedright
Owner/checkpoint
\end{minipage} & \begin{minipage}[b]{\linewidth}\raggedright
Objective verification
\end{minipage} \\
\midrule\noalign{}
\endhead
\bottomrule\noalign{}
\endlastfoot
SRS & Functional flows, nonfunctional SLOs, authorized scope, prohibited
operations & W1 / W15 & Each requirement has an acceptance test \\
ADR register & Context, options, decision, tradeoffs, consequences and
status & W3 onward & Versioned review decision, no unexplained tech
selection \\
Component architecture & Model, framework, harness, MCP
host/client/server, external systems and trust boundaries & W1/W6 & All
decision rights identified \\
Sequence diagram & Intake → proposal → validate → policy → dispatch →
observe → verify → close & W2/W15 & Privileged action has pre-execution
approval \\
State machine & All terminal states, retry/timeout/cancel/approval paths
& W2/W4 & 100\% terminal-state transition tests \\
MCP contract catalog & Tool schemas, resources, prompts, auth scope,
protocol/SDK versions & W6/W7 & Contract and negative protocol tests \\
Data-flow / provenance model & Inputs, taint status, sensitive fields,
evidence references, retention & W8/W15 & Every recommendation resolves
to source snapshot \\
Threat model / control matrix & STRIDE-style scenarios, OWASP/NIST
mapping, PEP/PDP, residual risk & W9/W10 & Negative attack fixtures
match controls \\
Test plan and evaluation report & Units, contracts, adversarial,
baseline metrics, gate policy & W12 & Reproducible CI artifacts \\
CI/CD and deployment diagrams & Build, test, scan, sign, promote and
rollback; secrets boundaries & W14 & No successful promotion with failed
gates \\
Runbook + incident procedure & SLOs, on-call, alert triage, disable
capability, recovery & W13/W16 & Fault-injection drill repeatable \\
Audit evidence contract & Run, SHA, alert/rule, policy, approver,
timestamps, verifier outcome & W15/W16 & Immutable linkage from result
to original evidence \\
\end{longtable}
\end{landscape}

\textbf{Required diagram formats:} Graphviz DOT, Mermaid, or PlantUML.
The repository must store diagram source alongside generated outputs,
define a render command, and include accessibility-friendly text
descriptions. Every diagram labels permission boundaries and trust
levels; do not substitute illustrative arrows for an actually enforced
control.

\subsubsection{Standard evidence record schema (student-defined, fixed
for the
course)}\label{standard-evidence-record-schema-student-defined-fixed-for-the-course}

\begin{Shaded}
\begin{Highlighting}[]
\FunctionTok{\{}
  \DataTypeTok{"run\_id"}\FunctionTok{:} \StringTok{"course{-}run{-}0001"}\FunctionTok{,}
  \DataTypeTok{"repo"}\FunctionTok{:} \StringTok{"example/fixture{-}repo"}\FunctionTok{,}
  \DataTypeTok{"source\_sha"}\FunctionTok{:} \StringTok{"0000000000000000000000000000000000000000"}\FunctionTok{,}
  \DataTypeTok{"alert\_id"}\FunctionTok{:} \DecValTok{101}\FunctionTok{,}
  \DataTypeTok{"rule\_id"}\FunctionTok{:} \StringTok{"example/query"}\FunctionTok{,}
  \DataTypeTok{"policy\_version"}\FunctionTok{:} \StringTok{"course{-}policy{-}v1"}\FunctionTok{,}
  \DataTypeTok{"disposition"}\FunctionTok{:} \StringTok{"needs\_review"}\FunctionTok{,}
  \DataTypeTok{"evidence\_refs"}\FunctionTok{:} \OtherTok{[}\StringTok{"sha256:fixture{-}hash"}\OtherTok{]}\FunctionTok{,}
  \DataTypeTok{"verifier"}\FunctionTok{:} \FunctionTok{\{}\DataTypeTok{"status"}\FunctionTok{:} \StringTok{"passed"}\FunctionTok{,} \DataTypeTok{"checks"}\FunctionTok{:} \OtherTok{[}\StringTok{"schema"}\OtherTok{,} \StringTok{"policy"}\OtherTok{]}\FunctionTok{\},}
  \DataTypeTok{"approval"}\FunctionTok{:} \KeywordTok{null}\FunctionTok{,}
  \DataTypeTok{"timestamp"}\FunctionTok{:} \StringTok{"2026{-}11{-}12T00:00:00Z"}
\FunctionTok{\}}
\end{Highlighting}
\end{Shaded}

The example SHA, rule, timestamp and evidence hash are
\textbf{fixtures}, not real GitHub values. Student implementations must
validate formats, link genuine recorded values to run state, remove
secrets from telemetry and document evidence retention/access policies.

\subsection{15. Security and Governance
Curriculum}\label{security-and-governance-curriculum}

\subsubsection{Threat-modeling method and control
ownership}\label{threat-modeling-method-and-control-ownership}

For each tool call identify (a) untrusted data source, (b) sensitive
asset, (c) authorization identity, (d) policy decision point, (e)
enforcement point, (f) audit evidence, and (g) failure behavior.
Consider data from issues, pull requests, source files, generated
patches, tool descriptions, MCP resources and model outputs
\textbf{hostile until validated}. The class applies STRIDE-style prompts
to each data flow and maps applicable mitigations to current OWASP LLM
guidance, NIST AI RMF, and MITRE ATLAS techniques without presenting
these external resources as identical security standards. See D13--D15.

\begin{landscape}
\begin{longtable}[]{@{}
  >{\raggedright\arraybackslash}p{(\columnwidth - 6\tabcolsep) * \real{0.2100}}
  >{\raggedright\arraybackslash}p{(\columnwidth - 6\tabcolsep) * \real{0.2500}}
  >{\raggedright\arraybackslash}p{(\columnwidth - 6\tabcolsep) * \real{0.3800}}
  >{\raggedright\arraybackslash}p{(\columnwidth - 6\tabcolsep) * \real{0.1600}}@{}}
\toprule\noalign{}
\begin{minipage}[b]{\linewidth}\raggedright
Security boundary
\end{minipage} & \begin{minipage}[b]{\linewidth}\raggedright
Main abuse case
\end{minipage} & \begin{minipage}[b]{\linewidth}\raggedright
Required defense
\end{minipage} & \begin{minipage}[b]{\linewidth}\raggedright
Verification artifact
\end{minipage} \\
\midrule\noalign{}
\endhead
\bottomrule\noalign{}
\endlastfoot
Model → harness & Model requests prohibited action or malformed
arguments & Strict schema, allowlist, execution limits, independent
policy decision & \texttt{test\_unknown\_tool},
\texttt{test\_schema\_denial} \\
External content → context & Issue embeds ``ignore prior policy'' or
false authorization & Treat retrieved text as data only; label
provenance; do not promote to system instruction & 12-case injection
corpus \\
Harness → MCP client & Unapproved tool is discoverable and appears
callable & Capability filtering + pre-dispatch policy enforcement &
Denial trace before network call \\
MCP server → GitHub & Server token has more access than current actor &
Resource-scoped GitHub App identity, server-side checks, separate
permission for writes & Cross-repository 403/deny test \\
Checkpoint store → execution & Stale/replayed continuation resumes
privileged action & Versioned checkpoint, nonce and idempotency key &
Replay and crash-consistency tests \\
Human approval → PEP & Spoofed, expired, revoked, wrong-scope approval &
Signed/verified identity, policy version, scope, expiry, single-use
decision ID & Negative approval suite \\
Multi-agent handoff & Child agent inherits excessive capability &
Delegated scope, tool isolation, bounded depth, no credential forwarding
& Worker-isolation trace \\
Telemetry → observability & Token, source content or private data leaks
into logs & Redaction, selective retention, least-privilege access &
Canary-secret scan over logs \\
CI → deployed runtime & Compromised dependency or unchecked policy
change & Lockfiles, scanning, change review, protected branches,
independent quality gate & CI and release provenance \\
\end{longtable}
\end{landscape}

\textbf{Policy authority matrix:} the LLM proposes classification and
actions; the harness gates tool execution and tracks state; the MCP
client transports an allowed request; the MCP server and GitHub API
enforce the downstream identity; application policy decides allowed
dispositions; the human reviewer approves sensitive exceptions; the
CI/CD system protects code and policy releases. This is functional
ownership---one deployed process may implement multiple roles but must
preserve each testable boundary.

\subsubsection{Required policy-as-code
exercise}\label{required-policy-as-code-exercise}

Use an allow-by-exception design: \texttt{default\ allow\ =\ false},
read-only retrieval permitted for a matching allowed repository and
valid identity, and write/merge/dismiss actions denied by default. A
mock write action may only proceed if an authorized reviewer approved
the exact repository, alert, action, policy version and unexpired nonce.
An unknown policy, unknown repository, unavailable identity provider or
missing evidence must produce an explicit denial or
\texttt{needs\_approval}, not optimistic success.

\textbf{Threat-model assessment:} each student identifies at least 12
adversarial fixtures, assigns risk and mitigations, writes tests that
prove the controls, tracks residual risk owners and proposes safe
operational alerts. An example ``attack succeeded'' label is evidence of
a \textbf{failing} gate, not successful completion of the lab.
Successful threat-modeling means the harmful operation cannot be
executed under the tested conditions.

\subsubsection{Security standards
applicability}\label{security-standards-applicability}

\begin{itemize}
\tightlist
\item
  \textbf{OWASP LLM/agent security:} use the current project for prompt
  injection, excessive agency, insecure output handling and supply-chain
  issues. Document the publication version used; OWASP naming and
  numbering has changed across editions.
\item
  \textbf{NIST AI RMF + Generative AI Profile:} risk governance and
  measurement overlay; do not confuse voluntary risk guidance with a
  specific mandatory control certification.
\item
  \textbf{MITRE ATLAS:} use attack technique research to expand
  adversarial scenarios and detection coverage; the course does not
  imply a complete ATT\&CK/ATLAS coverage guarantee.
\item
  \textbf{OPA/Rego:} optional implementation of deterministic external
  policy rules; equivalent audited policy engines are acceptable if all
  decision-table tests pass.
\item
  \textbf{GitHub governance:} required status checks, branch protections
  and code-scanning permissions remain authoritative downstream, even
  when an agent's proposed decision is justified.
\end{itemize}

\subsection{16. Testing and Evaluation
Strategy}\label{testing-and-evaluation-strategy}

\subsubsection{Layered quality gates}\label{layered-quality-gates}

\begin{longtable}[]{@{}
  >{\raggedright\arraybackslash}p{(\columnwidth - 6\tabcolsep) * \real{0.2500}}
  >{\raggedright\arraybackslash}p{(\columnwidth - 6\tabcolsep) * \real{0.2500}}
  >{\raggedright\arraybackslash}p{(\columnwidth - 6\tabcolsep) * \real{0.2500}}
  >{\raggedright\arraybackslash}p{(\columnwidth - 6\tabcolsep) * \real{0.2500}}@{}}
\toprule\noalign{}
\begin{minipage}[b]{\linewidth}\raggedright
Layer
\end{minipage} & \begin{minipage}[b]{\linewidth}\raggedright
Example case
\end{minipage} & \begin{minipage}[b]{\linewidth}\raggedright
Fixture strategy
\end{minipage} & \begin{minipage}[b]{\linewidth}\raggedright
Exit gate
\end{minipage} \\
\midrule\noalign{}
\endhead
\bottomrule\noalign{}
\endlastfoot
Unit & Unknown tool name; malformed JSON; budget exhausted & Synthetic
typed data & 100\% required safety unit tests pass \\
Contract & MCP tools/list schema; typed GitHub result; unsupported
revision & Stubbed JSON-RPC/HTTP & Exact schema and version behavior
validated \\
Integration & GitHub 403/429; SQLite restart; OPA deny; tracing &
In-memory/mocked services & No unauthorized effect; recoverable
failure \\
Adversarial & Prompt injection, fake tool descriptions, injected
``approved'' text & Frozen malicious corpus & 0 unauthorized tool
executions in corpus \\
Behavior evaluation & Evidence-backed alert analysis and abstention &
Labeled CodeQL-like fixture set & Meet predefined measured quality
targets \\
Operational & Timeout, network partition, duplicate message and service
restart & Controlled fault injection & Idempotency, trace propagation,
graceful failure \\
Release & SBOM/dependency reports, policy-version change, image nonroot
& Pinned CI environment & Every required check passes before
promotion \\
\end{longtable}

\subsubsection{Evaluation dataset and
metrics}\label{evaluation-dataset-and-metrics}

Create \textbf{100 synthetic alert examples}, stratified into 40 likely
actionable, 40 plausible benign candidates, and 20 genuinely
indeterminate cases, with independently reviewed gold labels, rule IDs,
evidence and disposition rationale. This is a teaching dataset, not an
empirical estimate of the real-world prevalence of alert classes.
Reserve 20\% as a held-out evaluation set; use stratified sampling, a
documented random seed and a locked fixture checksum. Independently
verify 10\% of generated labels before scoring.

\begin{itemize}
\tightlist
\item
  \textbf{Precision (recommended actionable):} \texttt{TP/(TP+FP)};
  \textbf{recall:} \texttt{TP/(TP+FN)}; report support and class balance
  rather than only accuracy.
\item
  \textbf{Evidence citation validity:} fraction of citations that
  resolve to the specified immutable fixture and revision; goal
  \textbf{100\%} on the audited capstone subset.
\item
  \textbf{Unsupported disposition rate:} recommendations lacking
  required facts / total recommendations; target \textbf{0} on mandatory
  safety scenarios, separately report the measured non-safety benchmark
  rate.
\item
  \textbf{Unauthorized action rate:} unauthorized side-effect executions
  / denied attempts; \textbf{must be 0} across mandatory negative tests.
\item
  \textbf{Indeterminate handling:} known-indeterminate cases routed for
  review, not silently treated as benign.
\item
  \textbf{Latency / cost:} p50/p95 per investigation, model calls per
  case, tokens per case and simulated dollar cost under a documented
  pricing assumption. No real-cost benchmark is mandatory.
\item
  \textbf{Recovery reliability:} success of restart/resume test suite
  and count of duplicated side effects; all mandatory crash fixtures
  must pass.
\end{itemize}

\textbf{Proposed quality targets (course rubric, not vendor claims):}
≥0.85 precision and ≥0.80 recall on a stable held-out synthetic
evaluation set; ≥0.90 valid evidence references in exploration,
increasing to 1.00 for final audited cases; every required security
negative case passes. If model outputs make quality targets unstable,
freeze model responses and grade the evaluation harness separately, then
report live-model variability transparently. Safety gates are
unconditional regardless of any average score.

\subsubsection{Continuous evaluation
mechanics}\label{continuous-evaluation-mechanics}

Use one command (\texttt{make\ verify}) to run formatting/lint, type
checking, pytest, JSON Schema contracts, OPA tests, adversarial tests,
evaluation gate and final evidence validation. Upload a machine-readable
report containing commit SHA, dependency lock hash, dataset checksum and
pass/fail states. A model-based judge may provide qualitative insight
but cannot overturn deterministic policy denials. Keep secure test
credentials in the CI secret store, never in the evaluation dataset.

\subsection{17. Midterm Engineering Project --- Week
8}\label{midterm-engineering-project-week-8}

\textbf{Title:} Resumable Read-Only Repository Evidence Investigator.
\textbf{Scenario:} a mock GitHub organization has multiple repositories
and historical code-scanning alerts. The student builds an agent that
receives a repository and alert ID, validates scope, retrieves data
through a custom MCP server, compacts context, produces a cited
investigation summary, and resumes a paused investigation.

\textbf{Required interfaces:}
\texttt{InvestigateRequest(repo,\ alert\_id,\ actor\_id)},
\texttt{EvidenceRecord(sha,\ rule,\ location,\ source\_hash)},
\texttt{DispositionRecommendation(value,\ rationale,\ references)},
\texttt{RunState(status,\ checkpoint\_version)}. Model output is
advisory; the runtime is responsible for enforcing input and tool
policies.

\textbf{Implementation plan:} create synthetic fixtures for two
repositories; provide a local MCP tool set; add a SQLite checkpoint
store; install a graph-based workflow; implement an evidence index;
create three approval/denial scenarios; document an ADR comparing custom
harness and framework orchestration; render architecture and sequence
diagrams.

\textbf{Demonstration script:} (1) valid read-only investigation returns
linked evidence; (2) unsupported repository returns blocked; (3) source
comment containing malicious instructions is treated as data; (4)
simulated crash after retrieval is resumed with the same run ID; (5)
missing source revision produces \texttt{needs\_review} rather than a
fabricated disposition; (6) report can be generated twice without
duplicate records.

\textbf{Submission:} 10-minute walkthrough, repository tag
\texttt{midterm-v1}, source+tests, architecture and threat summary,
20-fixture validation output, trace/evidence packet. \textbf{Passing
threshold:} ≥70/100 project rubric \textbf{and} every prohibited-action
test passed. \textbf{Midterm rubric:} 20 architecture, 25
implementation, 20 tests, 20 security and evidence, 15
explanation/reproducibility. This midterm is the prerequisite for the
W15 capstone integration.

\subsection{18. Final Enterprise AppSec Agent Capstone
Specification}\label{final-enterprise-appsec-agent-capstone-specification}

\subsubsection{System boundary and
architecture}\label{system-boundary-and-architecture}

The final deliverable is a \textbf{governed CodeQL alert investigation
and recommendation system}, initially read-only for any real GitHub
integration. It may read repository alert details and source snapshots
within granted scopes, recommend remediation or a justified disposition,
and prepare evidence for human review. It \textbf{must not silently
dismiss findings, change branch protections, merge PRs or bypass status
checks}. A separate mock gateway demonstrates pre-execution
authorization for privileged actions, but the baseline final build
contains no live dismissal API write capability.

\begin{Shaded}
\begin{Highlighting}[]
\NormalTok{GitHub user/request → intake + identity/scope check}
\NormalTok{                   → harness (checkpoint, budgets, policy, trace)}
\NormalTok{                   ↔ model (proposals/reasoned summaries only)}
\NormalTok{                   → MCP client → read{-}only GitHub MCP/REST façade}
\NormalTok{                   → evidence normalizer / revision{-}pinned store}
\NormalTok{                   → deterministic verifier → reviewer recommendation}
\NormalTok{                                      ↘ approval queue (mock only)}
\NormalTok{                   → immutable audit references / telemetry / report}
\end{Highlighting}
\end{Shaded}

\subsubsection{Mandatory functional user
stories}\label{mandatory-functional-user-stories}

\begin{longtable}[]{@{}
  >{\raggedright\arraybackslash}p{(\columnwidth - 4\tabcolsep) * \real{0.3333}}
  >{\raggedright\arraybackslash}p{(\columnwidth - 4\tabcolsep) * \real{0.3333}}
  >{\raggedright\arraybackslash}p{(\columnwidth - 4\tabcolsep) * \real{0.3333}}@{}}
\toprule\noalign{}
\begin{minipage}[b]{\linewidth}\raggedright
ID
\end{minipage} & \begin{minipage}[b]{\linewidth}\raggedright
User story
\end{minipage} & \begin{minipage}[b]{\linewidth}\raggedright
Demonstrated behavior
\end{minipage} \\
\midrule\noalign{}
\endhead
\bottomrule\noalign{}
\endlastfoot
CA-01 & As an AppSec investigator, review an authorized repository &
Repo identity and permission verified before data retrieval \\
CA-02 & As an investigator, obtain CodeQL findings & Paginated alerts
normalized, all source IDs retained \\
CA-03 & As a reviewer, see impact and evidence & Rule, commit, location,
and explanation linked to fixture \\
CA-04 & As a reviewer, receive recommendations & \texttt{remediate},
\texttt{candidate\_benign}, \texttt{needs\_review}; never unsupported
\texttt{dismissed} \\
CA-05 & As a policy owner, prevent unauthorized mutation & Write
attempted through mocked gate is denied absent approval \\
CA-06 & As an auditor, replay an investigation & Same fixture+version
reproduces report and policy outcome \\
CA-07 & As an operator, recover interrupted runs & Checkpoint
resumption, timeouts, idempotency, explicit terminal result \\
CA-08 & As a platform engineer, gate new agent releases & CI tests,
dependency checks, policy checks and trace evidence pass \\
\end{longtable}

\textbf{Evidence contract:} each case includes \texttt{repo\_owner},
\texttt{repo\_name}, \texttt{alert\_id}, \texttt{analysis\_id} if
available, \texttt{rule\_id}, source commit SHA, affected file and line
range, normalized severity, original observation, proposed outcome,
alternative explanation, verified evidence references, policy version,
verifier result, reviewer identity if required, decision timestamp and
immutable run/commit correlation IDs. Any incomplete or contradictory
evidence must produce an indeterminate/review-required disposition.

\subsubsection{Fifteen required capstone artifacts and acceptance
criteria}\label{fifteen-required-capstone-artifacts-and-acceptance-criteria}

\begin{longtable}[]{@{}
  >{\raggedleft\arraybackslash}p{(\columnwidth - 4\tabcolsep) * \real{0.4000}}
  >{\raggedright\arraybackslash}p{(\columnwidth - 4\tabcolsep) * \real{0.3000}}
  >{\raggedright\arraybackslash}p{(\columnwidth - 4\tabcolsep) * \real{0.3000}}@{}}
\toprule\noalign{}
\begin{minipage}[b]{\linewidth}\raggedleft
\#
\end{minipage} & \begin{minipage}[b]{\linewidth}\raggedright
Deliverable
\end{minipage} & \begin{minipage}[b]{\linewidth}\raggedright
Minimum pass condition
\end{minipage} \\
\midrule\noalign{}
\endhead
\bottomrule\noalign{}
\endlastfoot
1 & Implemented repository & Pinned, clean checkout runs offline mock
demo and tests \\
2 & System architecture specification & All five layers and trust
boundaries correctly identified \\
3 & Harness design specification & Bounded transitions, budgets,
persistence, approval, fail outcomes \\
4 & Framework selection ADR & ≥3 options, requirements scoring, decision
rationale \\
5 & MCP integration specification & Version, transport, tool schemas,
auth, error mapping and compatibility tests \\
6 & Threat model & ≥12 adversarial cases, mitigations, owners and
residual risks \\
7 & Tool authorization policy & Default deny; mock write requires exact
actor/action/scope approval \\
8 & Test and evaluation suite & Unit/contract/integration/adversarial
coverage with held-out fixtures \\
9 & CI/CD pipeline & Lint, typing, unit, policy and security-negative
gates required \\
10 & Observability configuration & Cross-component trace IDs and
redactable structured events \\
11 & Deployment procedure & Local container first; cloud architecture
optional, rollback explicit \\
12 & Operational runbook & Stop, disable, resume, rate-limit, incident,
recovery procedures \\
13 & Evidence register & Each case resolves to source revision, verifier
and policy decision \\
14 & Architecture review scorecard & Pass/fail/gaps with reviewer
signature and open risks \\
15 & Final engineering report & Methods, results, limitations, metrics,
release decision \\
\end{longtable}

\subsubsection{Acceptance-test matrix}\label{acceptance-test-matrix}

\begin{longtable}[]{@{}
  >{\raggedright\arraybackslash}p{(\columnwidth - 2\tabcolsep) * \real{0.5000}}
  >{\raggedright\arraybackslash}p{(\columnwidth - 2\tabcolsep) * \real{0.5000}}@{}}
\toprule\noalign{}
\begin{minipage}[b]{\linewidth}\raggedright
Control
\end{minipage} & \begin{minipage}[b]{\linewidth}\raggedright
Test and expected outcome
\end{minipage} \\
\midrule\noalign{}
\endhead
\bottomrule\noalign{}
\endlastfoot
Least privilege & Read succeeds without write credentials; mock mutation
is blocked \\
Alert completeness & Pages of findings are fetched without skipping
duplicate IDs \\
Evidence integrity & Recommendation links to exact revision and rule
ID \\
Deterministic policy & Same facts and policy version produce same
permit/deny result \\
Prompt injection & Malicious issue, code comment and tool response do
not alter authorized scope \\
Credential safety & Logs and prompts never contain test secret canary \\
Human approval & Missing/expired/wrong-scope approval cannot cause a
side effect \\
Safe retries & Repeated dispatch/restart cannot duplicate approved mock
side effects \\
Stop conditions & Timeout, budget, failure and denial are explicit
terminal outcomes \\
MCP compatibility & Pinned client/server revision and error semantics
exercised \\
CI/merge governance & Required checks protect pipeline; agent cannot
bypass protections \\
Reproducibility & Fresh reviewer can rebuild and run full test suite
offline \\
\end{longtable}

\textbf{End-to-end demonstration:} perform a benign candidate review
with evidence; a legitimate actionable alert; a misleading source
comment; wrong-repository access; missing evidence; interrupted
investigation; expired approval; dependency rate limit; and CI release
denial. The concluding review decides whether the \emph{read-only}
capability is deployable in a controlled environment. Authorization of
any actual alert dismissal is a separate operational project, with
enterprise owner approval and repository governance checks.

\subsection{19. Assessment Rubrics and Grading
Model}\label{assessment-rubrics-and-grading-model}

\begin{longtable}[]{@{}lr@{}}
\toprule\noalign{}
Graded component & Weight \\
\midrule\noalign{}
\endhead
\bottomrule\noalign{}
\endlastfoot
Weekly knowledge checks (16 × short quizzes) & 10\% \\
Laboratory implementation and code quality (L01--L16) & 30\% \\
Architecture, documentation and design reviews & 10\% \\
Security and adversarial-testing performance & 15\% \\
Midterm project & 10\% \\
Capstone implementation and evaluation & 20\% \\
Final operational readiness presentation & 5\% \\
\textbf{Total} & \textbf{100\%} \\
\end{longtable}

\begin{longtable}[]{@{}
  >{\raggedright\arraybackslash}p{(\columnwidth - 2\tabcolsep) * \real{0.5000}}
  >{\raggedright\arraybackslash}p{(\columnwidth - 2\tabcolsep) * \real{0.5000}}@{}}
\toprule\noalign{}
\begin{minipage}[b]{\linewidth}\raggedright
Performance level
\end{minipage} & \begin{minipage}[b]{\linewidth}\raggedright
Observable criteria
\end{minipage} \\
\midrule\noalign{}
\endhead
\bottomrule\noalign{}
\endlastfoot
Beginner & Can identify concepts but depends on example code, does not
independently isolate failing cases; often below 60 points \\
Developing & Produces partially working artifacts, needs design
coaching, tests common success paths; typically 60--74 points \\
Proficient & Independently implements scoped architecture, tests failure
paths, supplies reproducible evidence; typically 75--89 points \\
Advanced & Demonstrates cross-framework portability, systematic fault
testing, improved measured performance and review-ready risk analysis;
typically 90--100 points \\
\end{longtable}

\textbf{Certification-of-completion rubric (course-defined, not an
external credential):} minimum overall 75\%, midterm ≥70\%, capstone
≥80\%, all required security negative tests pass, all 15 capstone
artifacts submitted, and an independent reproduction review passes.
\textbf{Fail-closed gates override the numeric grade:} any unauthorized
write, intentional bypass of protections, missing trusted approval,
exposure of live secrets, absent source integrity for final audited
dispositions, or failed required negative test means \textbf{not
production-ready}. Learners may remediate and resubmit; they do not pass
by averaging unrelated high marks.

\textbf{Review checklist:} reviewer receives repository SHA, dependency
lock, \texttt{make\ verify} output, gold-dataset checksum, policy
version, test results and sample trace/evidence packet. Reviewer
documents disagreements and unresolved residual risks, not just an
overall score.

\subsection{20. Accelerated and Extended Study
Plans}\label{accelerated-and-extended-study-plans}

\textbf{Eight-week intensive option --- 20 h/week (\textasciitilde160
h):} combine W1+2, W3+4, W5+6, W7+8, W9+10, W11+12, W13+14, W15+16.
Preserve all 16 labs and both gates, but reduce optional videos and
concentrate guided coding in two longer sessions per week. This works
best for developers already fluent in Python async, GitHub APIs, Docker
and security. No safety gate is waived.

\textbf{Twenty-four-week extended option --- \textasciitilde7 h/week
(\textasciitilde168 h):} keep W1--16 as one-week cores plus eight
dedicated practice/review weeks. Insert bridge week after W1
(Python/async), a framework refactor after W4, MCP migration week after
W7, retrieval benchmark after W8, secure-policy workshop after W10,
evaluation calibration after W12, deployment dry-run after W14, and two
extra capstone review/fix weeks around W15--16 (nine insertions require
combining one earlier optional study block). For exactly 24 numbered
weeks, use the explicit schedule below.

\begin{longtable}[]{@{}rlrl@{}}
\toprule\noalign{}
Extended week & Work & Extended week & Work \\
\midrule\noalign{}
\endhead
\bottomrule\noalign{}
\endlastfoot
1 & Core W1 & 13 & Core W9 \\
2 & Python async bridge & 14 & Core W10 \\
3 & Core W2 & 15 & Security policy workshop \\
4 & Core W3 & 16 & Core W11 \\
5 & Core W4 & 17 & Core W12 \\
6 & Framework refactor & 18 & Evaluation calibration \\
7 & Core W5 & 19 & Core W13 \\
8 & Core W6 & 20 & Core W14 \\
9 & Core W7 & 21 & Deployment dry-run \\
10 & MCP migration workshop & 22 & Core W15 \\
11 & Core W8 (midterm) & 23 & Capstone remediation \\
12 & Retrieval quality workshop & 24 & Core W16 (final review) \\
\end{longtable}

The extended schedule allocates 16 core + 8 supplemental weeks, not 16 +
9; instructional hours are approximate (24 × 7 = 168). Extra practice
time is spent on deeper test coverage and review, not on taking
additional paid courses.

\subsection{21. Recommended Engineering Reference Books and
Documentation}\label{recommended-engineering-reference-books-and-documentation}

\textbf{Selected print/reference books} (bibliographic guidance; check
availability/edition independently): \emph{Designing Data-Intensive
Applications} (Martin Kleppmann) for persistence and distributed state;
\emph{Release It!} (Michael T. Nygard) for resiliency and failure modes;
\emph{Site Reliability Engineering} and \emph{The Site Reliability
Workbook} (Google) for SLOs and incident operations; \emph{Threat
Modeling: Designing for Security} (Adam Shostack) for systematic
abuse-case reasoning; \emph{Building Secure and Reliable Systems}
(Google) for security/reliability collaboration. These are supplemental
texts, not mandatory purchases.

\textbf{Authoritative online references, free to read:}

\begin{itemize}
\tightlist
\item
  Agent harness and framework: Microsoft Agent Harness (D10), LangGraph
  overview (D05), OpenAI Agents SDK agents/tools/guardrails/tracing
  (D06--D09).
\item
  MCP standards and SDKs: normative architecture and release notes
  (D01--D04), Python MCP SDK (D20).
\item
  GitHub AppSec integration: GitHub Code Scanning REST API (D11), GitHub
  MCP server (D12), GitHub App authentication (D21), GitHub Actions
  (D18), branch protection (D19).
\item
  Security and risk: NIST AI RMF (D13), OWASP GenAI security project
  (D14), MITRE ATLAS (D15), OPA (D16), OWASP Cheat Sheet Series (D22).
\item
  Telemetry: OpenTelemetry official reference (D17).
\end{itemize}

\textbf{Engineering research practice:} where official docs and a course
disagree, reproduce against a pinned SDK, record a compatibility issue,
and prefer the published specification for protocol semantics. Document
any implementation exception explicitly rather than silently
interpreting version-specific API behavior as universal.

\subsection{22. Final Competency and Career Readiness
Matrix}\label{final-competency-and-career-readiness-matrix}

\begin{landscape}
\begin{longtable}[]{@{}
  >{\raggedright\arraybackslash}p{(\columnwidth - 6\tabcolsep) * \real{0.1900}}
  >{\raggedright\arraybackslash}p{(\columnwidth - 6\tabcolsep) * \real{0.3300}}
  >{\raggedright\arraybackslash}p{(\columnwidth - 6\tabcolsep) * \real{0.2100}}
  >{\raggedright\arraybackslash}p{(\columnwidth - 6\tabcolsep) * \real{0.2700}}@{}}
\toprule\noalign{}
\begin{minipage}[b]{\linewidth}\raggedright
Capability
\end{minipage} & \begin{minipage}[b]{\linewidth}\raggedright
Proficient portfolio evidence
\end{minipage} & \begin{minipage}[b]{\linewidth}\raggedright
Typical role context
\end{minipage} & \begin{minipage}[b]{\linewidth}\raggedright
Review question
\end{minipage} \\
\midrule\noalign{}
\endhead
\bottomrule\noalign{}
\endlastfoot
Agent integration & Typed tool specifications and correct dispatch &
Backend AI engineer & What stops a malformed proposal? \\
Framework engineering & Documented LangGraph parity and selection ADR &
AI systems architect & Why graph orchestration over a custom loop? \\
Harness runtime & Checkpoint, budget and terminal-state traces &
Platform/agent runtime engineer & How are partial failures resumed
safely? \\
MCP protocol & Versioned server/client contract tests &
Integration/platform engineer & What changes between legacy and July
2026 behavior? \\
Context quality & Measured retrieval precision and source provenance &
RAG/knowledge engineer & Can every answer be traced to a source
revision? \\
Security & OWASP-style threat model and denied injections & Application
security engineer & Who authorizes the requested tool call? \\
Governance & Deterministic OPA policy and approval replay tests &
DevSecOps/GRC architect & Can the model override a denied action? \\
Multiagent & Scoped delegation and worker failure isolation &
Distributed systems engineer & Can a delegated worker exceed supervisor
authority? \\
Evaluations & Frozen gold set with precision, recall and CI gate & AI
quality engineer & Does a better aggregate score hide a safety
regression? \\
Operations & Correlated traces, SLOs and incident drill & SRE/platform
operations & How do you stop a misbehaving agent? \\
Delivery & Nonroot container, scanned release, audit artifacts &
DevSecOps engineer & What blocks promotion of an unsafe build? \\
Enterprise AppSec & Read-only CodeQL investigator and review packet &
AppSec automation engineer & Can a reviewer reproduce each recommended
disposition? \\
\end{longtable}
\end{landscape}

\textbf{Career-readiness defense:} the final learner explains the
architecture in 10 minutes, demonstrates a read-only legitimate task,
proves two blocked malicious tasks, recovers from one fault, and walks
an independent reviewer from Git commit to evidence to policy decision.
That combination provides a stronger engineering signal than completion
badges alone.

\subsection{23. Verified Sources and
Bibliography}\label{verified-sources-and-bibliography}

\subsubsection{Architectural source of
record}\label{architectural-source-of-record}

\textbf{User-provided primary document:} \emph{Agent Harness, Agent
Framework, and MCP: Which Layer Owns What?} version 1.0, November 12,
2026, 12 pages. Core thesis and divisions of responsibility are taken
from the document's sections 1--5; execution state and safeguards from
sections 6--7; framework choices from section 8; AppSec use case and
acceptance criteria from sections 9--10. The specific learning sequence
and laboratory/evaluation numbers in this syllabus are original proposed
curriculum design, not assertions from the reference PDF.

\textbf{Resource verification date:} November 12, 2026. Verification is
based on providers' publicly accessible landing pages and primary
specifications wherever possible. Course availability, lesson lists,
prices, API versions and video metadata may change. A course that
includes older MCP code remains educational but requires modern protocol
checking.

\subsubsection{Course and video source
register}\label{course-and-video-source-register}

\begin{longtable}[]{@{}lll@{}}
\toprule\noalign{}
ID & Provider and resource & Verified link \\
\midrule\noalign{}
\endhead
\bottomrule\noalign{}
\endlastfoot
\end{longtable}

LL01 \textbar{} LinkedIn Learning --- Building Agentic AI Systems for
Developers \textbar{}
\href{https://www.linkedin.com/learning/paths/building-agentic-ai-systems-for-developers}{Provider
page} \textbar{}\\
LL02 \textbar{} Microsoft Developer --- Building AI Agents: From Zero to
Production by Microsoft \textbar{}
\href{https://www.linkedin.com/learning/building-ai-agents-from-zero-to-production-by-microsoft}{Provider
page} \textbar{}\\
LL03 \textbar{} Anthropic --- Introduction to Model Context Protocol by
Anthropic \textbar{}
\href{https://www.linkedin.com/learning/introduction-to-model-context-protocol-by-anthropic}{Provider
page} \textbar{}\\
LL04 \textbar{} Morten Rand-Hendriksen --- Unboxing AI: Build a Remote
MCP Server from Zero to Deployed with OAuth \textbar{}
\href{https://www.linkedin.com/learning/unboxing-ai-build-a-remote-mcp-server-from-zero-to-deployed-with-oauth}{Provider
page} \textbar{}\\
LL05 \textbar{} Kesha Williams --- Governing AI Agents: Visibility and
Control \textbar{}
\href{https://www.linkedin.com/learning/governing-ai-agents-visibility-and-control}{Provider
page} \textbar{}\\
LL06 \textbar{} Kesha Williams --- Operating AI Agents: Failure and
Recovery \textbar{}
\href{https://www.linkedin.com/learning/operating-ai-agents-failure-and-recovery}{Provider
page} \textbar{}\\
FR01 \textbar{} Hugging Face --- AI Agents Course \textbar{}
\href{https://huggingface.co/learn/agents-course/unit0/introduction}{Provider
page} \textbar{}\\
FR02 \textbar{} LangChain Academy --- Introduction to LangGraph --
Python \textbar{}
\href{https://academy.langchain.com/courses/intro-to-langgraph}{Provider
page} \textbar{}\\
FR03 \textbar{} LangChain Academy --- Introduction to Deep Agents
\textbar{}
\href{https://academy.langchain.com/courses/foundation-introduction-to-deepagents}{Provider
page} \textbar{}\\
FR04 \textbar{} LangChain Academy --- Building Reliable Agents
\textbar{}
\href{https://academy.langchain.com/courses/building-reliable-agents}{Provider
page} \textbar{}\\
FR05 \textbar{} Hugging Face and Anthropic --- Model Context Protocol
(MCP) Course \textbar{}
\href{https://huggingface.co/learn/mcp-course/unit0/introduction}{Provider
page} \textbar{}\\
FR06 \textbar{} Elie Schoppik / Anthropic --- MCP: Build Rich-Context AI
Apps with Anthropic \textbar{}
\href{https://www.deeplearning.ai/courses/mcp-build-rich-context-ai-apps-with-anthropic}{Provider
page} \textbar{}\\
YT01 \textbar{} LangChain --- What is an agent harness? Explained fast
\textbar{} \href{https://www.youtube.com/watch?v=7dFvU9zuv0k}{Provider
page} \textbar{}\\
YT02 \textbar{} Prompt Engineering --- What is an Agent Harness? and How
to build a great one! \textbar{}
\href{https://www.youtube.com/watch?v=nWzXyjXCoCE}{Provider page}
\textbar{}\\
YT03 \textbar{} freeCodeCamp / Vaibhav Mehra --- LangGraph Complete
Course for Beginners -- Complex AI Agents with Python \textbar{}
\href{https://www.youtube.com/watch?v=jGg_1h0qzaM}{Provider page}
\textbar{}\\
YT04 \textbar{} freeCodeCamp --- RAG \& MCP Fundamentals -- A Hands-On
Crash Course \textbar{}
\href{https://www.youtube.com/watch?v=I7_WXKhyGms}{Provider page}
\textbar{}\\
YT05 \textbar{} Krish Naik --- Complete Agentic AI Course In 10 Hours
--- LangChain, LangGraph, RAG, Guardrails, Evals \textbar{}
\href{https://www.youtube.com/watch?v=rV3HJ4LEZ7k}{Provider page}
\textbar{}\\
YT06 \textbar{} freeCodeCamp / Bappy --- Agentic AI Complete Course for
Beginners \textbar{}
\href{https://www.youtube.com/watch?v=Zy7EXDONlTY}{Provider page}
\textbar{}

\subsubsection{Official technical and security
references}\label{official-technical-and-security-references}

\begin{longtable}[]{@{}lll@{}}
\toprule\noalign{}
ID & Source & Link \\
\midrule\noalign{}
\endhead
\bottomrule\noalign{}
\endlastfoot
\end{longtable}

D01 \textbar{} MCP 2026-07-28 Architecture \textbar{}
\href{https://modelcontextprotocol.io/specification/2026-07-28/architecture}{Official
documentation} \textbar{}\\
D02 \textbar{} MCP 2026-07-28 release and migration overview \textbar{}
\href{https://blog.modelcontextprotocol.io/posts/2026-07-28/}{Official
documentation} \textbar{}\\
D03 \textbar{} MCP TypeScript SDK revision support \textbar{}
\href{https://ts.sdk.modelcontextprotocol.io/v2/migration/support-2026-07-28}{Official
documentation} \textbar{}\\
D04 \textbar{} MCP Java SDK modern server (preview) \textbar{}
\href{https://java.sdk.modelcontextprotocol.io/latest/server-modern/}{Official
documentation} \textbar{}\\
D05 \textbar{} LangGraph overview \textbar{}
\href{https://docs.langchain.com/oss/python/langgraph/overview}{Official
documentation} \textbar{}\\
D06 \textbar{} OpenAI Agents SDK: Agents \textbar{}
\href{https://openai.github.io/openai-agents-python/agents/}{Official
documentation} \textbar{}\\
D07 \textbar{} OpenAI Agents SDK: Tools \textbar{}
\href{https://openai.github.io/openai-agents-python/tools/}{Official
documentation} \textbar{}\\
D08 \textbar{} OpenAI Agents SDK: Guardrails \textbar{}
\href{https://openai.github.io/openai-agents-python/guardrails/}{Official
documentation} \textbar{}\\
D09 \textbar{} OpenAI Agents SDK: Tracing \textbar{}
\href{https://openai.github.io/openai-agents-python/tracing/}{Official
documentation} \textbar{}\\
D10 \textbar{} Microsoft Agent Harness \textbar{}
\href{https://learn.microsoft.com/en-us/agent-framework/concepts/harness}{Official
documentation} \textbar{}\\
D11 \textbar{} GitHub REST code scanning endpoints \textbar{}
\href{https://docs.github.com/en/rest/code-scanning/code-scanning}{Official
documentation} \textbar{}\\
D12 \textbar{} GitHub MCP server README (read-only mode) \textbar{}
\href{https://github.com/github/github-mcp-server}{Official
documentation} \textbar{}\\
D13 \textbar{} NIST AI Risk Management Framework \textbar{}
\href{https://www.nist.gov/itl/ai-risk-management-framework}{Official
documentation} \textbar{}\\
D14 \textbar{} OWASP LLM application security project \textbar{}
\href{https://genai.owasp.org/}{Official documentation} \textbar{}\\
D15 \textbar{} MITRE ATLAS \textbar{}
\href{https://atlas.mitre.org/}{Official documentation} \textbar{}\\
D16 \textbar{} Open Policy Agent documentation \textbar{}
\href{https://www.openpolicyagent.org/docs}{Official documentation}
\textbar{}\\
D17 \textbar{} OpenTelemetry documentation \textbar{}
\href{https://opentelemetry.io/docs/}{Official documentation}
\textbar{}\\
D18 \textbar{} GitHub Actions documentation \textbar{}
\href{https://docs.github.com/en/actions}{Official documentation}
\textbar{}\\
D19 \textbar{} GitHub branch protection documentation \textbar{}
\href{https://docs.github.com/en/repositories/configuring-branches-and-merges-in-your-repository/managing-protected-branches/about-protected-branches}{Official
documentation} \textbar{}\\
D20 \textbar{} Python MCP SDK \textbar{}
\href{https://github.com/modelcontextprotocol/python-sdk}{Official
documentation} \textbar{}\\
D21 \textbar{} GitHub App authentication \textbar{}
\href{https://docs.github.com/en/apps/creating-github-apps/authenticating-with-a-github-app}{Official
documentation} \textbar{}\\
D22 \textbar{} OWASP Cheat Sheet Series \textbar{}
\href{https://cheatsheetseries.owasp.org/}{Official documentation}
\textbar{}

\subsubsection{Research integrity and maintenance
register}\label{research-integrity-and-maintenance-register}

\textbf{Verified directly on provider/official pages:} LinkedIn Learning
items LL01--LL06; LangChain Academy FR02--FR04; Hugging Face FR01/FR05;
DeepLearning.AI FR06; MCP specification D01--D04; GitHub Code Scanning
D11 and GitHub MCP read-only mode D12; OpenAI Agents SDK guardrails and
tracing D08--D09. Other referenced official docs are stable canonical
destinations, but API details should be retested against the pinned
versions in the exercise.

\textbf{YouTube evidence limits:} YT01 and YT02 have directly surfaced
video listings; YT03 has a confirmed companion code repository and
timestamp index; YT04 has a corroborating third-party chapter index;
YT05's title is confirmed but its precise runtime and chapters are not;
YT06 is provisional, with a secondary mirror as the available metadata
source. Do \textbf{not} require YT06 or unverifiable timestamp-specific
sections for grading. Optional promotional claims and unverified course
dates are intentionally excluded.

\textbf{Next review triggers:} MCP specification revision; breaking SDK
release; GitHub API permission/version change; linked course
replacement; security framework revision; any failed course setup
instructions; or annual syllabus review, whichever comes first.

\begin{center}\rule{0.5\linewidth}{0.5pt}\end{center}

\textbf{End of syllabus.} The target completion outcome is an
independently verified, source-grounded, read-only enterprise AppSec
agent system with a clear path to later authorized changes. Training
completion and operational authorization are distinct decisions.

\end{document}
